% concatenated from Prots.arxiv.tex
\pdfoutput=1
\documentclass[11pt]{amsart}

\usepackage[T1]{fontenc}
\usepackage[utf8]{inputenc}
\usepackage[main=english]{babel}
\usepackage{lmodern}
\usepackage[a4paper,margin=1in]{geometry}
\usepackage[final,protrusion=true,expansion=false]{microtype}
\emergencystretch=2em

\usepackage{amsmath,amssymb,mathtools}
\usepackage{amsthm}
\usepackage{enumitem}
\usepackage[sort,compress]{cite}

\usepackage[hypertexnames=false]{hyperref}
\hypersetup{%
	colorlinks=true,
	linkcolor=blue,
	citecolor=blue,
	urlcolor=blue,
	pdftitle={A Bundle-Theoretic Formulation of Phonons in Crystalline Phases},
	pdfauthor={Aleksey Prots},
	pdfkeywords={crystalline phase, phonons, torus bundle, Goldstone modes, connection, jet bundle}
}

\theoremstyle{plain}
\newtheorem{proposition}{Proposition}[section]
\theoremstyle{definition}
\newtheorem{definition}[proposition]{Definition}
\theoremstyle{remark}
\newtheorem{remark}[proposition]{Remark}

\title[Bundle-Theoretic Formulation of Phonons]{A Bundle-Theoretic Formulation of Phonons in Crystalline Phases}
\author{Aleksey Prots}
\address{Department of Theoretical Physics and Computer Technologies, Kuban State University, Krasnodar, Russian Federation}
\email{BWPSM.lab@gmail.com}
\date{April 28, 2026}
\subjclass[2020]{53C05, 55R10, 74B05, 74J05}
\keywords{Crystalline phase, order parameter, phonons, Goldstone modes, frame-bundle reduction, associated torus bundle, Ehresmann connection, jet bundle, Lagrangian field theory}

\begin{document}

\begin{abstract}
	Phonons are usually introduced by choosing a local displacement field. 
	This paper keeps that local description, but identifies the global geometric 
	object represented by it. The aim is not to change the local acoustic 
	equations, but to describe the global configuration space of the translational 
	order parameter on a fixed crystallographic background and to give a globally 
	defined replacement for the displacement gradient.
	
	After the orientational part of the crystalline order has been fixed by a 
	reduction of the orthonormal frame bundle to a discrete point group, the 
	translational order parameter is described as a section of an associated 
	torus bundle. In a symmorphic crystal the point group acts on the translation 
	torus linearly, whereas in a nonsymmorphic crystal the action is affine and 
	records the extension class of the crystallographic group. Relative to the 
	fixed point-group bundle, the discreteness of the structure group gives a 
	canonical flat Ehresmann connection on the associated torus bundle. The 
	corresponding covariant differential of the translational field is a globally 
	defined object which locally coincides with the ordinary displacement gradient.
	
	This covariant differential is then used to formulate the phonon sector as a 
	first-order Lagrangian field theory. When the flat torus holonomy fixes an 
	equilibrium point, linearization about the corresponding covariantly constant 
	section gives the usual local displacement field. For derivative-only quadratic 
	elastic Lagrangians satisfying the standard objectivity condition, the theory 
	reduces locally to linear elasticity and to the standard acoustic phonon 
	spectrum. If such a global equilibrium section does not exist, the same 
	linear theory is understood locally on defect-free simply connected patches.
\end{abstract}



\maketitle

\section{Introduction}

Phonons are usually described locally by displacement fields whose derivatives enter the elastic energy. 
This local description is sufficient for deriving the standard equations of elasticity. 
The purpose of this paper is to determine the global geometric object represented by such a displacement field when the crystalline background has nontrivial orientational structure.

Crystalline order contains two distinct components. 
The first component is translational. The position of the crystal relative to its lattice is defined only modulo a lattice vector; hence the natural target is not \(\mathbb R^d\), but the torus
\[
T_\Pi=\mathbb{R}^d/\Pi,
\]
where \(\Pi\subset\mathbb{R}^d\) is the translation lattice. 
The second component is orientational. A crystalline phase does not preserve the full rotational symmetry, but only a discrete point group \(P\). 
Geometrically, this orientational order is described by a reduction of the orthonormal frame bundle to \(P\).

The full order parameter of a crystalline phase may therefore be written schematically as
\[
\Phi=(u,\sigma_P),
\]
where \(u\) is the translational part and \(\sigma_P\) encodes the orientational reduction. This paper does not try to study the full coupled dynamics of both components. Instead, the orientational sector is fixed once and for all. Equivalently, we choose a principal \(P\)-bundle \(Q\) and treat it as part of the background. The remaining dynamical variable is the translational order parameter, and globally it is a section of the associated torus bundle
\[
\pi_Y:Y_\Pi=Q\times_P T_\Pi\to X.
\]
In this sense the paper studies the phonon sector on a fixed crystallographic background. The variable denoted by \(u\) should be understood as the translational phase of the crystalline order, not as a full nonlinear deformation map of a body into Euclidean space.

\medskip

\noindent\textbf{Relation to previous work and novelty.}
The local theory recovered below is the standard harmonic theory of lattice dynamics and elastic waves in crystals \cite{BornHuang,LandauLifshitzElasticity,NyePhysicalProperties,MusgraveCrystalAcoustics}. 
The interpretation of phonons as Goldstone modes of broken translations, and the elimination of independent rotational Goldstone variables by inverse-Higgs-type relations, is also standard in effective descriptions of spacetime symmetry breaking \cite{LeutwylerPhonons,LowManohar,WatanabeMurayama2013,BraunerInverseHiggs,BraunerBook}. The broader nonlinear-realization framework goes back to the CCWZ construction \cite{ColemanWessZumino1969,CallanColemanWessZumino1969}, while modern effective-field-theory discussions of solids and elastic variables provide closely related formulations \cite{NicolisEtAl2015,HayataHidaka2014}. 
A separate geometric literature describes dislocations and disclinations by torsion, curvature, and non-Riemannian geometric structures \cite{BilbyBulloughSmith1955,Kroner1981,deWit1973,KatanaevVolovich1992,Katanaev2005,KleinertGaugeFields,MuraMicromechanics}. Recent dual descriptions relating elasticity and defects to fracton tensor gauge theories give another related but distinct perspective \cite{PretkoRadzihovsky2018,PretkoZhaiRadzihovsky2019,PretkoChenYou2020,GromovRadzihovsky2024}. 
The present paper occupies a more elementary but useful layer between these two viewpoints: the defect-free phonon sector is formulated globally before defects or dynamical orientational variables are introduced. 
The novelty is therefore not a new local phonon spectrum. 
It is the identification of the translational order parameter with a section of an associated torus bundle, including affine nonsymmorphic gluing, and the construction of the canonical covariant differential whose local representative is the ordinary displacement gradient.

The main results are the following.
\begin{enumerate}[label=\textup{(\roman*)}]
\item The global configuration space of the translational order parameter on the fixed orientational background is the associated torus bundle
\[
Y_\Pi=Q\times_P T_\Pi.
\]
\item The symmorphic and nonsymmorphic cases are distinguished by the action of \(P\) on \(T_\Pi\): the former gives a linear torus action, while the latter gives an affine torus action determined by the crystallographic extension.
\item Relative to the fixed \(P\)-bundle \(Q\), the discreteness of \(P\) gives a canonical flat Ehresmann connection \(\Gamma_0\) on \(Y_\Pi\to X\).
\item For every section \(\sigma:X\to Y_\Pi\), the covariant differential \(\nabla^{\Gamma_0}\sigma\) is a globally defined section of
\[
T^*X\otimes(Q\times_P\mathbb R^d),
\]
and in admissible local trivializations it is represented by the ordinary first derivative of the local displacement field.
\end{enumerate}

The construction does not change the local phonon equations once the usual elastic assumptions are imposed. After choosing an admissible local trivialization and linearizing about a covariantly constant equilibrium section, derivative-only quadratic Lagrangians satisfying the standard objectivity condition reproduce the usual displacement field, the small-strain tensor, the elastic equations, and the acoustic dispersion relations.
Its role is to identify the global configuration bundle for the translational order parameter and the covariant differential that gives the invariant meaning of the local displacement gradient.

Three points should be fixed at the outset. 
First, we assume that the orthonormal frame bundle admits a global reduction \(Q\subset F_g(M)\) to the point group \(P\). 
This is a nontrivial topological condition. If such a reduction does not exist, or if defects are present, the construction should be understood on the defect-free region. 
Second, in the nonsymmorphic case the point group acts on \(T_\Pi\) by affine rather than purely linear transformations. 
This affine action is retained in the definition of \(Y_\Pi\). 
Its differential is the usual linear point-group action on vertical tangent spaces. 
Consequently, the local linear phonon theory depends only on the linear part of the action, whereas the affine shifts enter the global gluing of the torus-valued field. 
Third, the global linearization around an undeformed reference section requires a covariantly constant equilibrium section. Such a section exists only when the holonomy of the flat torus bundle fixes a point of \(T_\Pi\). If this fixed-point condition fails, the linear phonon theory remains a local theory, or a theory on a chosen defect-free simply connected region with a chosen background lift.

The formulation uses the language of associated bundles, Ehresmann connections, and first jet bundles. 
In these terms the phonon sector becomes a first-order Lagrangian field theory on \(J^1Y_\Pi\). 
Physically, this gives a global geometric formulation of the standard Goldstone interpretation of phonons: they are the low-energy modes associated with broken translations. 
The rotational Goldstone variables are not independent acoustic modes; in the low-energy theory they are expressed through derivatives of the displacement field.

The structure of the paper is as follows. Sec.~\ref{sec:crystal_symmetry} fixes the crystallographic data and the distinction between symmorphic and nonsymmorphic cases. Sec.~\ref{sec:orderparameter} explains the translational and orientational parts of the crystalline order parameter. Sec.~\ref{sec:associated_bundle} constructs the associated torus bundle and the canonical flat connection. Sec.~\ref{sec:first_order_lagrangian} formulates the first-order Lagrangian theory and derives the local phonon limit. The cubic-crystal check of the local elastic limit is given in Appendix~\ref{app:cubic_example}.

Throughout the paper, \(M\) denotes the spatial base manifold on which the crystal is described, whereas \(X\) denotes the base of the field theory.
In the static case \(X=M\).
In the dynamical case,
\[
X=\mathbb{R}_t\times M.
\]
Spatial geometric structures over \(M\) are then pulled back to \(X\) along the projection
\[
\mathrm{pr}_M:\mathbb{R}_t\times M\to M.
\]

\section{Crystallographic Symmetry Data}
\label{sec:crystal_symmetry}

A crystalline phase is characterized by a residual discrete subgroup of the Euclidean symmetry group. 
In an idealized isotropic medium the full Euclidean group acts as a spatial symmetry, whereas in a crystal only the motions compatible with the lattice remain: translations by lattice vectors and finite rotations belonging to the point group. 
Crystallization may therefore be viewed as a standard example of spontaneous breaking of spatial symmetries: continuous translations are reduced to a lattice, and continuous rotations are reduced to a finite subgroup of \(O(d)\) \cite{Chaikin,BornHuang,LeutwylerPhonons,NicolisEtAl2015}.

\medskip

\subsection{Spatial setting}
\label{subsec:spatial_setting}
Let \(E\) be a Euclidean space of dimension \(d\), identified with \(\mathbb{R}^d\) after a choice of origin. 
The residual symmetry of an ideal crystal is described by a discrete subgroup
\[
\Gamma \subset \mathrm{E}(d)=\mathbb{R}^d\rtimes O(d),
\]
where \(\mathrm{E}(d)\) denotes the group of Euclidean motions of \(E\).

\begin{definition}[Crystallographic group]
	\label{def:crystallographic_group}
	A crystallographic group is a discrete subgroup
	\[
	\Gamma\subset \mathrm{E}(d)
	\]
	acting on the Euclidean space \(E\) by isometries such that the orbit space \(E/\Gamma\) is compact. 
	Equivalently, by the Bieberbach theorems, \(\Gamma\) contains a full-rank translation subgroup of finite index \cite{Charlap1986}.
\end{definition}

In particular, the subgroup of pure translations
\begin{equation}
	\Pi:=\Gamma\cap \mathbb{R}^d
	\label{eq:lattice_def}
\end{equation}
is a full-rank lattice, \(\Pi\simeq \mathbb{Z}^d\). 
The quotient
\begin{equation}
	P:=\Gamma/\Pi
	\label{eq:point_group_def}
\end{equation}
is finite and is realized, through the linear parts of the elements of \(\Gamma\), as a subgroup of \(O(d)\). 
It is called the point group of the crystal. 
Thus \(\Gamma\) fits into the short exact sequence
\begin{equation}
	1\longrightarrow \Pi \longrightarrow \Gamma \longrightarrow P \longrightarrow 1.
	\label{eq:cryst_extension}
\end{equation}
In this group-extension formulation, the symmorphic case corresponds to a split extension, so the crystallographic group has the form
\[
\Gamma \simeq \Pi \rtimes P.
\]
If the extension does not split, such a semidirect-product presentation is not available in general. 
This nonsymmorphic case is represented by nontrivial affine translational data; geometrically, it includes screw axes and glide planes \cite{BradleyCracknell,InternationalTablesA}. 
For the general theory of crystallographic and space groups, see, for example, \cite{Charlap1986,BradleyCracknell,InternationalTablesA}.

\subsection{Symmorphic and Nonsymmorphic Cases}
\label{subsec:symmorphic_nonsymmorphic_affine}

The short exact sequence \eqref{eq:cryst_extension} contains more information than the linear action of the point group on the lattice. 
This additional information is encoded by the extension class of \eqref{eq:cryst_extension}, or equivalently by a group \(2\)-cocycle. 
It is this class that distinguishes the symmorphic and nonsymmorphic cases. 
For the construction below, this distinction has to be realized not only algebraically, but also as an action on the translation torus \(T_\Pi=\mathbb R^d/\Pi\). 
We use the standard description of group extensions by cocycles; see, for example, \cite{Charlap1986,BrownCohomology}.

Write an element of the Euclidean group as
\[
(a,A)\in \mathbb R^d\rtimes O(d),
\]
where \(a\in\mathbb R^d\) is the translation part and \(A\in O(d)\) is the orthogonal part. 
It acts on \(x\in\mathbb R^d\) by
\[
x\longmapsto Ax+a,
\]
and the group law is
\[
(a,A)(b,B)=(a+Ab,AB).
\]
The inclusion \(\Pi\subset\Gamma\) identifies each lattice vector \(\lambda\in\Pi\) with the pure translation \((\lambda,I)\). 
For every coset \(p\in P=\Gamma/\Pi\), choose a representative
\[
\gamma_p=(a_p,A_p)\in\Gamma.
\]
The orthogonal part \(A_p\) is independent of the chosen representative, because changing the representative by a lattice translation changes only \(a_p\). 
Moreover, conjugation by \(\gamma_p\) sends lattice translations to lattice translations:
\[
\gamma_p(\lambda,I)\gamma_p^{-1}=(A_p\lambda,I),
\qquad
\lambda\in\Pi.
\]
Hence \(A_p\Pi=\Pi\). 
Thus the point group acts on the lattice by orthogonal automorphisms,
\[
A:P\longrightarrow \mathrm{Aut}(\Pi)\subset O(d),
\qquad
p\longmapsto A_p.
\]

\begin{proposition}[Affine action of the point group on the translation torus]
	\label{prop:affine_point_action}
	Let \(\Gamma\subset \mathbb R^d\rtimes O(d)\) be a crystallographic group. 
	Let
	\[
	\Pi=\Gamma\cap\mathbb R^d
	\]
	be its translation lattice, and let
	\[
	P=\Gamma/\Pi
	\]
	be its point group. 
	For each \(p\in P\), choose a representative
	\[
	\gamma_p=(a_p,A_p)\in\Gamma.
	\]
	Then the following statements hold.
	
	\begin{enumerate}[label=\textup{(\roman*)}]
		\item The map \(p\mapsto A_p\) defines a representation
		\[
		A:P\to \mathrm{Aut}(\Pi)\subset O(d).
		\]
		
		\item The element
		\begin{equation}
			c(p,q):=a_p+A_pa_q-a_{pq}
			\label{eq:extension_cocycle}
		\end{equation}
		belongs to \(\Pi\) and defines a \(2\)-cocycle of the group \(P\) with values in the \(P\)-module \(\Pi_A\), where the action of \(P\) on \(\Pi\) is given by the representation \(A\).
		
		\item Changing the representatives \(\gamma_p\) changes \(c\) by a coboundary. Hence the class
		\[
		[c]\in H^2(P,\Pi_A)
		\]
		is an invariant of the extension \eqref{eq:cryst_extension}. The extension splits, that is \(\Gamma\simeq\Pi\rtimes P\), if and only if this class is zero.
		
		\item The formula
		\begin{equation}
			\rho(p)[v]=[A_pv+a_p],
			\qquad [v]\in T_\Pi=\mathbb R^d/\Pi,
			\label{eq:affine_torus_action}
		\end{equation}
		defines a left affine action of \(P\) on the torus \(T_\Pi\). 
		In the symmorphic case, after a suitable choice of origin and representatives, one has \(a_p=0\), and \(\rho\) becomes the linear action \([v]\mapsto[A_pv]\).
	\end{enumerate}
\end{proposition}

\begin{proof}
	For \(\lambda\in\Pi\), conjugation by \(\gamma_p=(a_p,A_p)\) gives
	\[
	\gamma_p(\lambda,I)\gamma_p^{-1}=(A_p\lambda,I)\in\Pi.
	\]
	Hence \(A_p\Pi=\Pi\). Moreover, replacing \(\gamma_p\) by \((b_p,I)\gamma_p\), \(b_p\in\Pi\), changes only the translation part \(a_p\), not the orthogonal part \(A_p\). Thus \(p\mapsto A_p\) is well defined as a map \(P\to\mathrm{Aut}(\Pi)\subset O(d)\).
	
	For \(p,q\in P\), multiplication in \(\mathbb R^d\rtimes O(d)\) gives
	\[
	\gamma_p\gamma_q
	=
	(a_p+A_pa_q,A_pA_q).
	\]
	Since \(\gamma_p\gamma_q\) and \(\gamma_{pq}\) project to the same element \(pq\in P\), the element \(\gamma_p\gamma_q\gamma_{pq}^{-1}\) lies in \(\Pi\). Its orthogonal part is therefore the identity, so \(A_pA_q=A_{pq}\). Thus \(A:P\to\mathrm{Aut}(\Pi)\) is a representation. Its translation part gives
	\[
	c(p,q):=a_p+A_pa_q-a_{pq}\in\Pi.
	\]
	Associativity of multiplication in \(\Gamma\) gives
	\[
	A_pc(q,r)-c(pq,r)+c(p,qr)-c(p,q)=0,
	\]
	hence \(c\in Z^2(P,\Pi_A)\).
	
	If the representatives are changed to \(\gamma'_p=(b_p,I)\gamma_p\), \(b_p\in\Pi\), then \(a'_p=b_p+a_p\), and the corresponding cocycle is
	\[
	c'(p,q)
	=
	c(p,q)+b_p+A_pb_q-b_{pq}.
	\]
	Thus \(c'\) and \(c\) differ by a coboundary, so the class \([c]\in H^2(P,\Pi_A)\) is independent of the chosen representatives.
	
	The extension \eqref{eq:cryst_extension} splits if and only if \([c]=0\). Indeed, if \([c]=0\), the representatives may be changed so that \(c(p,q)=0\) for all \(p,q\). Then \(\gamma_p\gamma_q=\gamma_{pq}\), and the representatives form a subgroup of \(\Gamma\) isomorphic to \(P\). Conversely, if the extension splits, representatives can be chosen inside such a subgroup, and then \(c=0\).
	
	It remains to check the action on \(T_\Pi\). The formula
	\[
	\rho(p)[v]=[A_pv+a_p]
	\]
	is well defined with respect to \(v\): if \(v\) is replaced by \(v+\lambda\), \(\lambda\in\Pi\), then \(A_p\lambda\in\Pi\). It is also independent of the chosen representative \(\gamma_p\), since replacing \(a_p\) by \(a_p+b_p\), \(b_p\in\Pi\), does not change the class in \(\mathbb R^d/\Pi\). Finally,
	\[
	\begin{aligned}
		\rho(p)\rho(q)[v]
		&=[A_p(A_qv+a_q)+a_p]  \\
		&=[A_{pq}v+a_{pq}+c(p,q)] \\
		&=[A_{pq}v+a_{pq}]
		=\rho(pq)[v],
	\end{aligned}
	\]
	because \(c(p,q)\in\Pi\) vanishes in the quotient. Therefore \(\rho\) is a left affine action of \(P\) on \(T_\Pi\).
\end{proof}

\medskip

The preceding construction can be phrased cohomologically in a way that makes precise in what sense the affine torus action records the nonsymmorphic data. 
Let
\[
0\longrightarrow \Pi_A \longrightarrow \mathbb R^d_A \longrightarrow T_\Pi \longrightarrow 0
\]
be the exact sequence of \(P\)-modules, where \(P\) acts on \(\mathbb R^d\) and \(\Pi\) by the linear representation \(A\). Here \(T_\Pi\) is regarded as a \(P\)-module through the same linear action induced by \(A\); the affine translational parts then define a cocycle with values in this linear \(P\)-module.
The classes \(\bar a_p=[a_p]\in T_\Pi\) form a \(1\)-cocycle with values in \(T_\Pi\). 
Indeed, equation \eqref{eq:extension_cocycle} gives
\[
\bar a_{pq}=\bar a_p+A_p\bar a_q.
\]
The connecting homomorphism
\[
\delta:H^1(P,T_\Pi)\longrightarrow H^2(P,\Pi_A)
\]
associated with the above exact sequence sends the class of \(\bar a\) to the extension class \([c]\) represented by \eqref{eq:extension_cocycle}. 
Since \(P\) is finite and \(\mathbb R^d_A\) is a real \(P\)-module, averaging cochains gives
\[
H^n(P,\mathbb R^d_A)=0,
\qquad n>0.
\]
Consequently, in the relevant part of the long exact sequence, the connecting homomorphism identifies the affine translational class in \(H^1(P,T_\Pi)\) with the crystallographic extension class in \(H^2(P,\Pi_A)\). 
Thus the phrase ``the affine action records the extension class'' means that the translational part of the affine action determines, through this connecting map, the cohomology class distinguishing the nonsymmorphic extension from a split semidirect product.

\begin{remark}[Practical consequence for the translational order parameter]
	\label{rem:nonsymmorphic_affine_action}
	If only the linear action \([v]\mapsto[A_pv]\) is used below, then the theory is effectively restricted to the symmorphic case, or equivalently it forgets the extension class \([c]\). 
	For a general crystallographic group, the action of \(P\) on \(T_\Pi\) must be understood as the affine action \eqref{eq:affine_torus_action}. 
	Its differential is the linear action \(A_p\). 
	Thus the vertical tangent bundle, the covariant differential, and the linear phonon spectrum depend on \(A_p\), whereas the nonsymmorphic shifts \(a_p\) enter the gluing laws for the local representatives of the torus-valued field. 
	This distinction is built into the definition of \(Y_\Pi=Q\times_P T_\Pi\).
\end{remark}

\medskip

\noindent\textbf{Example: screw-axis gluing.}
Let the lattice contain a primitive vector \(e_z\) along an axis, and let \(R\) be a rotation of order \(n\) about this axis. 
A screw generator of type \(n_m\) may be represented by
\begin{equation}
	\gamma_s=\left(\frac{m}{n}e_z,R\right),
	\qquad
	R^n=I,
	\label{eq:screw_generator}
\end{equation}
with \(m e_z\in\Pi\). Then
\begin{equation}
	\gamma_s^{\,n}=(m e_z,I)\in\Pi,
	\label{eq:screw_power_translation}
\end{equation}
so the image \(s\in P\) has order \(n\), while its representative in \(\Gamma\) contains a fractional translation. The induced action on the translation torus is
\begin{equation}
	\rho(s)[v]=\left[Rv+\frac{m}{n}e_z\right],
	\qquad
	\rho(s)^n=\mathrm{id}_{T_\Pi}.
	\label{eq:screw_affine_torus_action}
\end{equation}
Thus, on an overlap whose transition function is \(s\), a local lift of the torus-valued displacement changes by the rotation \(R\), by the fractional shift \((m/n)e_z\), and by a lattice vector. 
Differentiating removes the constant fractional shift and leaves only the linear action of \(R\) on vertical tangent vectors. 
This is the local mechanism by which nonsymmorphic data affect the global gluing of \(\sigma\), while the principal symbol of the linearized phonon equations depends only on the linear part.

\medskip

\noindent\textbf{Scope of the main construction.}
The construction below uses only the translation lattice \(\Pi\), the point group \(P\), and a fixed principal \(P\)-bundle \(Q\). 
Magnetic or time-reversing symmetries would require replacing \(P\) by the appropriate magnetic point group, and are not considered here.
\medskip

\section{The Order Parameter as a Geometric Object}
\label{sec:orderparameter}

Geometrically, an order parameter associated with the breaking of a structure group \(G\) to a closed subgroup \(H\) is described by a section of a quotient bundle. 
Let
\[
\pi_{\mathcal P} : \mathcal P \to X
\]
be a principal \(G\)-bundle over \(X\).
The right action of \(H\) on \(\mathcal P\) defines the quotient bundle
\begin{equation}
	\pi_H:\mathcal P/H\longrightarrow X,
	\label{eq:quotient_bundle}
\end{equation}
whose typical fiber is the homogeneous space \(G/H\). 
A global section
\begin{equation}
	\sigma:X\longrightarrow \mathcal P/H
	\label{eq:general_order_parameter}
\end{equation}
is then interpreted as an order-parameter field. 
Equivalently, such a section determines a reduction of the structure group from \(G\) to \(H\), that is, a principal \(H\)-subbundle \(\mathcal P_H\subset\mathcal P\). 
This is the standard bundle-theoretic form of symmetry reduction used in the geometric description of Higgs fields and spontaneous symmetry breaking \cite{HusemollerFibreBundles,KobayashiNomizu1,SardanashvilyBook,GiachettaBook,BraunerBook}.

\medskip

For crystals this general scheme has to be adapted. 
The continuous group of spatial motions of the Euclidean space \(E\simeq \mathbb R^d\) is
\begin{equation}
	G=\mathrm E(d)=\mathbb R^d\rtimes O(d),
	\label{eq:euclidean_group}
\end{equation}
whereas the residual symmetry of an ideal crystal is a discrete crystallographic group
\[
\Gamma\subset \mathrm E(d),
\]
related to the translation lattice \(\Pi\) and the point group \(P\) by the exact sequence \eqref{eq:cryst_extension}. 
Thus, in the present setting, crystalline order is best separated into two geometrically distinct parts: the translational Goldstone sector and the orientational reduction of the frame bundle.

Locally, this separation leads to the following model of the space of crystalline order. 
The translational component is described by a point of the torus \(\mathbb R^d/\Pi\), while the orientation of the elementary cell, modulo the point-group action, is described by a point of \(O(d)/P\). 
Thus the local state space has the form
\begin{equation}
	\mathcal O_{\mathrm{cryst}}
	\simeq
	(\mathbb R^d/\Pi)\times (O(d)/P).
	\label{eq:order_space}
\end{equation}
This is only a local model of crystalline order. 
More invariantly, before this separation one may regard the ideal order-parameter space as the homogeneous space \(\mathrm E(d)/\Gamma\). The product expression in \eqref{eq:order_space} is obtained after choosing local representatives and local trivializations. In the nonsymmorphic case it should not be read as a global product decomposition: the translational torus is glued over changes of crystalline frame by the affine \(P\)-action induced by the crystallographic extension. 
Globally, the translational and orientational components are realized by different geometric objects.

\medskip

Let \(X\) denote the base of the theory, as in the introduction. 
On an open set \(U\subset X\), a local crystalline order parameter is therefore represented by a pair
\begin{equation}
	(u,\varrho),
	\qquad
	u:U\to \mathbb R^d/\Pi,
	\qquad
	\varrho:U\to O(d)/P.
	\label{eq:local_order_parameter}
\end{equation}
Here \(u\) is the local translational component, while \(\varrho\) is the local orientational component.

The translational component \(u\) is a displacement field defined modulo lattice translations. 
Locally, over an open set \(U\subset X\), it may be viewed as a section of the trivial torus bundle
\begin{equation}
	U\times (\mathbb R^d/\Pi)\longrightarrow U.
	\label{eq:trivial_torus_bundle}
\end{equation}
Physically, \(u\) describes translational order. 
It is the Goldstone field associated with the spontaneous breaking of the continuous translation group \(\mathbb R^d\) to the lattice \(\Pi\), and its small fluctuations are the phonon degrees of freedom \cite{LeutwylerPhonons,WatanabeMurayama2013,AndersenBraunerHofmannVuorinen,HayataHidaka2014}. 
Although crystallization also breaks continuous rotations, the rotational Goldstone variables do not give independent acoustic modes. 
In the low-energy theory they are expressed through derivatives of the displacement field, in accordance with the inverse Higgs mechanism and the redundancy of some spacetime Goldstone variables \cite{IvanovOgievetsky1975,LowManohar,WatanabeMurayama2013,BraunerInverseHiggs}. 
Globally, however, the translational component is not naturally a section of the trivial bundle \eqref{eq:trivial_torus_bundle}; its local representatives are glued by the point-group action determined by the orientational structure.

\medskip

We now specify the geometric meaning of the orientational component. 
Let a Riemannian metric \(g\) be fixed on the spatial base manifold \(M\). 
Then the orthonormal frames form a principal bundle
\begin{equation}
	\pi_F:F_g(M)\longrightarrow M
	\label{eq:orthonormal_frame_bundle}
\end{equation}
with structure group \(O(d)\). 
If an orientation is fixed, one may instead use the oriented frame bundle
\begin{equation}
	F_g^+(M)\longrightarrow M
	\label{eq:oriented_frame_bundle}
\end{equation}
with structure group \(SO(d)\). 
In the dynamical setting, these spatial bundles are pulled back to \(X=\mathbb R_t\times M\) along the projection \(\mathrm{pr}_M:X\to M\); for simplicity, we keep the same notation for the pullbacks.

Since \(O(d)\) acts on \(O(d)/P\) by left multiplication, the orientational order is described by the associated bundle
\begin{equation}
	\mathcal Y_P
	:=
	F_g(M)\times_{O(d)}\bigl(O(d)/P\bigr)
	\longrightarrow M,
	\label{eq:YP_bundle}
\end{equation}
where
\begin{equation}
	F_g(M)\times_{O(d)}\bigl(O(d)/P\bigr)
	:=
	\bigl(F_g(M)\times (O(d)/P)\bigr)\big/\sim,
	\qquad
	(f\cdot a,\,[b])\sim (f,\,[ab]),
	\quad a,b\in O(d).
	\label{eq:associated_equiv_relation}
\end{equation}
The fiber of \(\mathcal Y_P\) over \(x\in M\) is the space of possible orientations of the crystalline cell at \(x\), modulo the point group \(P\). 
Thus the orientational component of crystalline order is not an abstract map into \(O(d)/P\), but a section
\begin{equation}
	\sigma_P:M\longrightarrow \mathcal Y_P
	=
	F_g(M)\times_{O(d)}\bigl(O(d)/P\bigr).
	\label{eq:sigmaP}
\end{equation}
In the dynamical setting \(X=\mathbb R_t\times M\), the same object may be viewed either as a time-dependent family of sections of \(\mathcal Y_P\to M\), or equivalently as a section of the pullback bundle \(\mathrm{pr}_M^*\mathcal Y_P\to X\).

By the standard correspondence between reductions of structure group and sections of the associated quotient bundle, a section \(\sigma_P\) of \(\mathcal Y_P\to M\) is equivalent to a reduction of the \(O(d)\)-structure of \(F_g(M)\) to the subgroup \(P\). 
Equivalently, it determines a principal \(P\)-subbundle
\begin{equation}
	Q\subset F_g(M).
	\label{eq:Q_in_frame_bundle}
\end{equation}
Thus the orientational order may be described either by the section \(\sigma_P\) or by the reduced frame bundle \(Q\).

The global existence of such a reduction is a nontrivial topological condition. 
Concretely, it means that, after choosing crystalline frames, the transition functions of the orthonormal frame bundle can be taken to have values in the finite subgroup \(P\). 
This is not an automatic property of an arbitrary Riemannian manifold, and in the presence of defects the construction should be understood on the defect-free region. 
In this paper \(Q\) is treated as part of the fixed crystallographic background.

\medskip

Thus, in the crystalline case, the general apparatus of symmetry reduction remains useful, but the order parameter naturally separates into two components. 
Locally, these components may be written as
\begin{equation}
	u:U\to \mathbb R^d/\Pi,
	\qquad
	\sigma_P:M\to F_g(M)\times_{O(d)}\bigl(O(d)/P\bigr).
	\label{eq:two_components_order_parameter}
\end{equation}
Here \(u\) denotes a local representative of the translational component, while \(\sigma_P\) describes the orientational reduction. 
Schematically, the crystalline order parameter is therefore
\begin{equation}
	\Phi=(u,\sigma_P).
	\label{eq:order_parameter_pair_local}
\end{equation}

\begin{remark}
	\label{rem:scope_main_text}
	The full crystalline order parameter has two components,
	\[
	\Phi=(u,\sigma_P),
	\]
	where \(u\) represents translational order and \(\sigma_P\) represents orientational order. 
	In the rest of the paper the orientational component is fixed. 
	Equivalently, the principal \(P\)-bundle \(Q\subset F_g(M)\) is treated as part of the background data. 
	The remaining dynamical variable is the translational component, globally represented by a section
	\[
	\sigma:X\to Y_\Pi,
	\qquad
	Y_\Pi=Q\times_P T_\Pi.
	\]
	This section is the configuration field of the phonon sector.
\end{remark}

This is the decomposition used below. 
The orientational reduction \(Q\subset F_g(M)\) is fixed as background data, and the translational component is the field to be varied. 
The bundle \(Q\) supplies the twisting data for the global configuration space of this field.

\section{The Associated Bundle for a Discrete Structure Group}
\label{sec:associated_bundle}

From now on, the orientational sector is fixed. 
Equivalently, the principal \(P\)-bundle
\[
Q\subset F_g(M)
\]
is treated as part of the background data. 
Locally, the translational component can be represented by a displacement field with values in the fixed torus
\[
T_\Pi=\mathbb R^d/\Pi.
\]
Globally, this description is generally insufficient. 
The local representatives of the displacement field are defined relative to crystalline frames, and under a change of such a frame they are transformed by the point-group action of \(P\).

The section \(\sigma_P\) determines the principal \(P\)-subbundle
\[
Q\subset F_g(M),
\]
which records the allowed local crystalline frames. 
The point group acts on \(T_\Pi=\mathbb R^d/\Pi\): linearly in the symmorphic case and affinely in the nonsymmorphic case, as explained in Remark~\ref{rem:nonsymmorphic_affine_action}. 
Therefore the translational component is not, in general, a map into a single fixed torus. 
It is a section of the bundle associated to \(Q\).

\medskip

In the dynamical setting \(X=\mathbb R_t\times M\), the principal \(P\)-bundle over \(X\) is the pullback
\[
\mathrm{pr}_M^*Q\longrightarrow X,
\qquad
\mathrm{pr}_M:X=\mathbb R_t\times M\to M.
\]
For brevity, this pullback will again be denoted by \(Q\to X\). 
Let
\begin{equation}
	\rho:P\to \mathrm{Diff}(T_\Pi)
	\label{eq:rho_action}
\end{equation}
be the action of the point group on the translation torus. 
In the symmorphic case, \(\rho(p)[v]=[A_pv]\). 
In the nonsymmorphic case, \(\rho(p)[v]=[A_pv+a_p]\). 
The differential of \(\rho(p)\) is the linear map \(A_p\); hence vertical tangent directions are acted on only by the linear part.

\begin{definition}
	\label{def:associated_torus_bundle}
	The bundle associated to the principal \(P\)-bundle \(Q\to X\) with typical fiber \(T_\Pi\) is
	\begin{equation}
		Y_\Pi:=Q\times_P T_\Pi:=(Q\times T_\Pi)/P,
		\label{eq:Ypi}
	\end{equation}
	where
	\begin{equation}
		(q,\tau)\sim (qp,\rho(p^{-1})\tau),
		\qquad p\in P.
		\label{eq:Ypi_equiv}
	\end{equation}
	It is a fiber bundle over \(X\) with projection
	\begin{equation}
		\pi_Y:Y_\Pi\to X,
		\qquad
		\pi_Y([q,\tau])=\pi_Q(q).
		\label{eq:Ypi_projection}
	\end{equation}
\end{definition}

A configuration of the translational sector is a section of this bundle:
\[
\sigma:X\to Y_\Pi.
\]
In local trivializations it is represented by torus-valued displacement functions, but globally these local representatives are glued by the \(P\)-action.

Let \(\{U_\alpha\}\) be an open cover of \(X\) with local sections
\[
z_\alpha:U_\alpha\to Q.
\]
Each \(z_\alpha\) gives a trivialization
\begin{equation}
	\psi_\alpha:\pi_Y^{-1}(U_\alpha)\xrightarrow{\sim} U_\alpha\times T_\Pi
	\label{eq:local_trivialization_map}
\end{equation}
defined by
\begin{equation}
	\psi_\alpha([z_\alpha(x),\tau])=(x,\tau).
	\label{eq:local_trivialization_formula}
\end{equation}
On an overlap \(U_\alpha\cap U_\beta\), write
\[
z_\beta(x)=z_\alpha(x)g_{\alpha\beta}(x),
\qquad
g_{\alpha\beta}:U_\alpha\cap U_\beta\to P.
\]
Then the corresponding fiber coordinates are related by
\begin{equation}
	\tau_\beta
	=
	\rho(g_{\alpha\beta}(x)^{-1})\,\tau_\alpha .
	\label{eq:transition_functions_associated}
\end{equation}
Since \(P\) is discrete, the transition functions \(g_{\alpha\beta}\) are locally constant. Hence \(Y_\Pi\to X\) has locally constant transition functions.

A section of \(Y_\Pi\to X\) is a map
\begin{equation}
	\sigma:X\to Y_\Pi,
	\qquad
	\pi_Y\circ \sigma=\mathrm{id}_X.
	\label{eq:section_Ypi}
\end{equation}
In the local trivializations \eqref{eq:local_trivialization_map}, it is represented by maps
\[
\sigma_\alpha:U_\alpha\to T_\Pi
\]
satisfying the transition rule \eqref{eq:transition_functions_associated}. 
If \(\widetilde u_\alpha:U_\alpha\to \mathbb R^d\) is a local lift of \(\sigma_\alpha\), then on overlaps
\begin{equation}
	\widetilde u_\beta
	=
	A_{g_{\alpha\beta}}^{-1}\widetilde u_\alpha
	+
	a(g_{\alpha\beta}^{-1})
	+
	\lambda_{\alpha\beta},
	\qquad
	\lambda_{\alpha\beta}\in \Pi.
	\label{eq:local_lift_relation}
\end{equation}
Here \(A_{g_{\alpha\beta}}\) is the orthogonal part of the transition function, and \(a(g_{\alpha\beta}^{-1})\) is the affine shift appearing in the action of \(g_{\alpha\beta}^{-1}\) on \(T_\Pi\). 
In the symmorphic case this shift is zero. 
Thus local displacement lifts are glued by a point-group transformation, possibly by a nonsymmorphic shift, and by an arbitrary lattice vector.

\medskip

We next record the differential-geometric structure of \(Y_\Pi\). 
Since \(Y_\Pi\) is obtained as the quotient of \(Q\times T_\Pi\) by the action
\[
(q,\tau)\cdot p=(qp,\rho(p^{-1})\tau),
\]
its tangent bundle may be written as
\[
TY_\Pi\cong (TQ\times TT_\Pi)/P.
\]
Since \(T_\Pi=\mathbb R^d/\Pi\) is a torus, its tangent bundle is canonically trivial:
\[
TT_\Pi\cong T_\Pi\times \mathbb R^d.
\]
Consequently,
\begin{equation}
	TY_\Pi \cong (TQ\times T_\Pi\times \mathbb R^d)/P.
	\label{eq:TYpi_final}
\end{equation}

Consider the pullback of the tangent bundle of the base,
\begin{equation}
	\pi_Y^*TX:=Y_\Pi\times_X TX.
	\label{eq:pullback_TX}
\end{equation}
The differential of the projection \(\pi_Y\) defines a vector-bundle morphism over \(Y_\Pi\),
\begin{equation}
	T\pi_Y:TY_\Pi\longrightarrow \pi_Y^*TX,
	\qquad
	T\pi_Y(v_y)=(y,T_y\pi_Y(v_y)).
	\label{eq:TpiY}
\end{equation}
The vertical tangent bundle is its kernel:
\begin{equation}
	VY_\Pi:=\ker(T\pi_Y)\subset TY_\Pi.
	\label{eq:vertical}
\end{equation}
Since \(\pi_Y:Y_\Pi\to X\) is a fiber bundle, \(T\pi_Y\) is surjective. 
Thus one obtains the short exact sequence
\begin{equation}
	0\longrightarrow VY_\Pi\longrightarrow TY_\Pi
	\xrightarrow{\,T\pi_Y\,}\pi_Y^*TX\longrightarrow 0.
	\label{eq:exact}
\end{equation}

In the present case the vertical bundle admits a natural associated description:
\begin{equation}
	VY_\Pi \cong (Q\times TT_\Pi)/P.
	\label{eq:vertical_associated}
\end{equation}
Since \(TT_\Pi\) is canonically trivial for a torus, vertical directions are modeled by the vector space \(\mathbb R^d\), with the point-group action given by the differential of \(\rho\). 
In the nonsymmorphic case this differential removes the affine shifts; only the linear part \(A_p\) acts on vertical tangent vectors.

Vertical directions in \(VY_\Pi\) are tangent to the fibers \(T_\Pi\) and represent infinitesimal translational variations of the crystalline order. 
After linearization around a background section, these variations give the phonon degrees of freedom associated with broken continuous translations \cite{LeutwylerPhonons,WatanabeMurayama2013,HayataHidaka2014,AndersenBraunerHofmannVuorinen}.

The exact sequence \eqref{eq:exact} shows that a differential-geometric description of the phonon sector requires a horizontal subbundle, equivalently a splitting of this sequence. 
In the present case such a splitting is canonical: since the structure group \(P\) is discrete, the transition functions \eqref{eq:transition_functions_associated} are locally constant. 
The next subsection shows that this gives a canonical flat Ehresmann connection on \(\pi_Y:Y_\Pi\to X\).

\subsection{Horizontal Splitting and the Ehresmann Connection}
\label{subsec:ehresmann}

Consider the associated bundle
\[
\pi_Y:Y_\Pi\to X,
\qquad
Y_\Pi=Q\times_P T_\Pi,
\]
constructed above. 
For a general fiber bundle, the short exact sequence \eqref{eq:exact} has no canonical splitting. 
In the present case, however, the discreteness of \(P\) determines a natural horizontal distribution on \(Y_\Pi\to X\).

Let \(\{U_\alpha\}\) be an open cover of \(X\) such that \(Y_\Pi\) is trivialized over each \(U_\alpha\) by \eqref{eq:local_trivialization_map}. 
In such a trivialization one has the natural splitting
\begin{equation}
	T(U_\alpha\times T_\Pi)\cong TU_\alpha \oplus TT_\Pi.
	\label{eq:local_splitting}
\end{equation}
We define the local horizontal subbundle by
\begin{equation}
	H_\alpha:= TU_\alpha \oplus 0 \subset T(U_\alpha\times T_\Pi).
	\label{eq:Halpha}
\end{equation}

We now check that these subbundles are compatible on overlaps. 
The transition functions of \(Y_\Pi=Q\times_P T_\Pi\) are
\begin{equation}
	\Phi_{\alpha\beta}(x,\tau)
	=
	\bigl(x,\rho(g_{\alpha\beta}(x)^{-1})\tau\bigr),
	\label{eq:Phi_transition}
\end{equation}
where
\[
g_{\alpha\beta}:U_\alpha\cap U_\beta\to P
\]
are the transition functions of the principal \(P\)-bundle \(Q\to X\). 
Since \(P\) is discrete, \(g_{\alpha\beta}\) is locally constant. 
Thus, on each connected component of \(U_\alpha\cap U_\beta\),
\begin{equation}
	\Phi_{\alpha\beta}(x,\tau)
	=
	(x,\varphi_{\alpha\beta}(\tau)),
	\qquad
	\varphi_{\alpha\beta}:=\rho(g_{\alpha\beta}^{-1}),
	\label{eq:varphi_alphabeta}
\end{equation}
where \(\varphi_{\alpha\beta}\in \mathrm{Diff}(T_\Pi)\) is independent of \(x\).
Therefore the differential of the transition map is
\begin{equation}
	T\Phi_{\alpha\beta}(v,\xi)
	=
	\bigl(v,\,T_\tau\varphi_{\alpha\beta}(\xi)\bigr),
	\qquad
	(v,\xi)\in TU_\alpha\oplus TT_\Pi.
	\label{eq:TPhi}
\end{equation}
In particular,
\begin{equation}
	T\Phi_{\alpha\beta}(v,0)=(v,0).
	\label{eq:horizontal_invariance}
\end{equation}
Thus the subspaces \(TU_\alpha\oplus 0\) are preserved by the transition functions. 
Consequently, the local subbundles \(H_\alpha\) agree on overlaps and glue to a global subbundle
\begin{equation}
	HY_\Pi \subset TY_\Pi.
	\label{eq:horizontal}
\end{equation}

Locally, under the trivialization 
\(\pi_Y^{-1}(U_\alpha)\simeq U_\alpha\times T_\Pi\), the vertical bundle is
\begin{equation}
	VY_\Pi|_{\pi_Y^{-1}(U_\alpha)}
	\cong
	0\oplus TT_\Pi.
	\label{eq:vertical_local}
\end{equation}
Thus, in each trivialization,
\begin{equation}
	T(U_\alpha\times T_\Pi)
	\cong
	(TU_\alpha\oplus 0)\oplus (0\oplus TT_\Pi).
	\label{eq:local_direct_sum}
\end{equation}
After gluing, this gives the global decomposition
\begin{equation}
	TY_\Pi = HY_\Pi \oplus VY_\Pi.
	\label{eq:splitting}
\end{equation}

\begin{proposition}[Canonical flat connection]
	\label{prop:ehresmann}
	Let
	\[
	Y_\Pi=Q\times_P T_\Pi
	\]
	be the bundle associated to a principal \(P\)-bundle \(Q\to X\), where \(P\) is discrete. 
	Then the local horizontal subbundles \eqref{eq:Halpha} agree on overlaps and define a global subbundle
	\[
	HY_\Pi\subset TY_\Pi
	\]
	such that
	\[
	TY_\Pi=HY_\Pi\oplus VY_\Pi.
	\]
	Moreover, \(HY_\Pi\) is integrable. 
	Thus, relative to the fixed \(P\)-bundle \(Q\), it defines a canonical flat Ehresmann connection on \(\pi_Y:Y_\Pi\to X\).
\end{proposition}

\begin{proof}
	In a local trivialization \(Y_\Pi|_{U_\alpha}\simeq U_\alpha\times T_\Pi\), set \(H_\alpha=TU_\alpha\oplus 0\). 
	The transition functions \(\Phi_{\alpha\beta}\) act as the identity on the base coordinates and by locally constant diffeomorphisms \(\varphi_{\alpha\beta}\) on the fiber coordinates. 
	Hence their differentials preserve the subspaces \(TU_\alpha\oplus 0\). 
	Thus the subbundles \(H_\alpha\) glue to a global subbundle \(HY_\Pi\subset TY_\Pi\). 
	The local decomposition \eqref{eq:local_direct_sum} then gives the global decomposition \eqref{eq:splitting}.
	
	To verify integrability, observe that in each local trivialization the horizontal vector fields have the form \((X,0)\), where \(X\) is a vector field on \(U_\alpha\). 
	For two such fields,
	\[
	[(X,0),(Y,0)] = ([X,Y],0).
	\]
	Thus \(H_\alpha\) is involutive in every trivialization, and the glued distribution \(HY_\Pi\) is integrable. 
	Therefore \(HY_\Pi\) defines a flat Ehresmann connection on \(\pi_Y:Y_\Pi\to X\).
\end{proof}

The word ``canonical'' is used relative to the fixed \(P\)-bundle \(Q\). 
It should not be read as a canonical choice before the crystallographic background \(Q\) has been specified, nor as a choice canonical simultaneously for all possible \(P\)-reductions. 
In admissible local trivializations the transition functions are the identity on the base and are locally constant in the fiber variable. 
Hence the subspaces \(TU_\alpha\oplus 0\) glue without any additional choice. 
This flat Ehresmann connection is associated with the discrete \(P\)-bundle \(Q\); it is not the Levi--Civita connection of the Riemannian frame bundle \(F_g(M)\). 
Relating it to the Levi--Civita connection would require additional assumptions ensuring that Levi--Civita parallel transport preserves the \(P\)-reduction \(Q\).

Restricting \(T\pi_Y\) to the horizontal subbundle gives an isomorphism
\begin{equation}
	T\pi_Y|_{HY_\Pi}:HY_\Pi\xrightarrow{\sim}\pi_Y^*TX,
	\label{eq:TpiY_horizontal_iso}
\end{equation}
which is equivalent to a splitting of the short exact sequence \eqref{eq:exact}.
We now rewrite this connection in the jet-bundle language used for first-order Lagrangians. 
In this language, an Ehresmann connection on \(Y_\Pi\to X\) is equivalently a global section of
\[
J^1Y_\Pi\to Y_\Pi .
\]
This is the form used below to define the covariant differential and to write Lagrangians depending on first derivatives.

\subsection{The Covariant Differential in the Jet Formulation}
\label{subsec:jet_covariant}

The horizontal distribution constructed in Sec.~\ref{subsec:ehresmann},
\[
HY_\Pi \subset TY_\Pi,
\]
defines a canonical flat Ehresmann connection on \(\pi_Y:Y_\Pi\to X\). 
We now rewrite this connection in the jet-bundle language used for first-order Lagrangians. 
In this language, a connection on a bundle \(Y\to X\) is a global section of the affine bundle of first jets
\[
J^1Y \to Y;
\]
see, for example, \cite{SardanashvilyBook,GiachettaBook}.

Let
\[
\pi_Y:Y_\Pi\to X
\]
be the configuration bundle, and denote its first jet bundle by \(J^1Y_\Pi\). 
For \(y\in Y_\Pi\) with \(x=\pi_Y(y)\), the fiber \((J^1Y_\Pi)_y\) consists of \(1\)-jets of local sections \(\sigma\) satisfying \(\sigma(x)=y\). 
Thus \(J^1Y_\Pi\) records both the value of the field and its first derivatives.

The bundle
\[
J^1Y_\Pi \to Y_\Pi
\]
is an affine bundle modeled on the vector bundle
\begin{equation}
	\pi_Y^*T^*X \otimes_{Y_\Pi} VY_\Pi \to Y_\Pi.
	\label{eq:jet_affine_model}
\end{equation}
In other words, the difference of two \(1\)-jets over the same point \(y\in Y_\Pi\) is a vertical-valued covector.

\medskip

In this terminology, a connection on \(Y_\Pi\to X\) is a global section
\begin{equation}
	\Gamma_Y: Y_\Pi \to J^1Y_\Pi.
	\label{eq:connection_jet}
\end{equation}
Equivalently, it is a horizontal distribution \(HY_\Pi\subset TY_\Pi\), or a splitting of the short exact sequence \eqref{eq:exact}.

In a local trivialization
\[
Y_\Pi|_{U_\alpha}\simeq U_\alpha\times T_\Pi
\]
choose coordinates \((x^\lambda,\tau^i)\), where \(x^\lambda\) are coordinates on \(U_\alpha\subset X\) and \(\tau^i\) are local coordinates on \(T_\Pi\). 
The first jet bundle \(J^1Y_\Pi\) then carries induced coordinates
\[
(x^\lambda,\tau^i,\tau^i_\lambda),
\]
where \(\tau^i_\lambda\) denote the jet coordinates corresponding to the first derivatives \(\partial_\lambda\tau^i\).

In these coordinates the connection \eqref{eq:connection_jet} is written as
\begin{equation}
	\Gamma_Y
	=
	(x^\lambda,\tau^i,\Gamma^i_\lambda(x,\tau)).
	\label{eq:connection_coordinates}
\end{equation}
The corresponding covariant differential is the first-order operator
\begin{equation}
	D_{\Gamma_Y}: J^1Y_\Pi \to \pi_Y^*T^*X \otimes VY_\Pi,
	\label{eq:DGamma}
\end{equation}
given locally by
\begin{equation}
	D_{\Gamma_Y}
	=
	(\tau^i_\lambda-\Gamma^i_\lambda)\,
	dx^\lambda\otimes \partial_i.
	\label{eq:DGamma_formula}
\end{equation}
It measures the difference between an arbitrary \(1\)-jet and the \(1\)-jet selected by the connection.

\medskip

Here \(u\) denotes a local physical displacement representative, while \(\sigma\) denotes a global section of \(Y_\Pi\to X\). 
Let
\[
\sigma:X\to Y_\Pi
\]
be such a section. 
Its first jet
\begin{equation}
	J^1\sigma:X\to J^1Y_\Pi
	\label{eq:first_jet_sigma}
\end{equation}
defines the covariant differential
\begin{equation}
	\nabla^\Gamma\sigma
	:=
	D_\Gamma\circ J^1\sigma.
	\label{eq:covariant}
\end{equation}
In local coordinates this is
\begin{equation}
	\nabla^\Gamma\sigma
	=
	(\partial_\lambda \sigma^i
	-
	\Gamma^i_\lambda\circ \sigma)\,
	dx^\lambda\otimes \partial_i.
	\label{eq:covariant_local}
\end{equation}
Thus \(\nabla^\Gamma\sigma\) is the global counterpart of the local displacement gradient; it is defined independently of a chosen trivialization of \(Y_\Pi\).

\medskip

The case of principal interest is the canonical flat Ehresmann connection constructed in Sec.~\ref{subsec:ehresmann}. 
Since the transition functions of \(Y_\Pi=Q\times_P T_\Pi\) are locally constant, the horizontal distribution in the local trivializations \(U_\alpha\times T_\Pi\) is
\[
H_\alpha = TU_\alpha\oplus 0.
\]
The corresponding jet connection \(\Gamma_0:Y_\Pi\to J^1Y_\Pi\) has the local form
\begin{equation}
	\Gamma_0
	=
	(x^\lambda,\tau^i,0),
	\label{eq:flat_connection_coordinates}
\end{equation}
or equivalently
\[
\Gamma_{0\,\lambda}^i=0.
\]
This expression is compatible on overlaps because the fiber transition functions are locally constant in \(x\). Hence
\begin{equation}
	D_{\Gamma_0}
	=
	\tau^i_\lambda\,dx^\lambda\otimes \partial_i,
	\qquad
	\nabla^{\Gamma_0}\sigma
	=
	\partial_\lambda \sigma^i\,
	dx^\lambda\otimes \partial_i.
	\label{eq:flat_covariant_differential}
\end{equation}
Thus, for the canonical flat connection, the covariant differential is locally the ordinary first derivative with respect to the base coordinates. 
Its significance is that this local expression is induced by a global connection on \(Y_\Pi\to X\), and therefore transforms correctly on overlaps.

\medskip

Since \(\sigma\) is a section of \(Y_\Pi\to X\), the covariant differential \(\nabla^\Gamma\sigma\) is a vertical-valued one-form along \(\sigma\). 
Set
\begin{equation}
	\sigma^*VY_\Pi
	:=
	X\times_{Y_\Pi}VY_\Pi.
	\label{eq:pullback_vertical}
\end{equation}
This vector bundle over \(X\) has fibers
\begin{equation}
	(\sigma^*VY_\Pi)_x
	=
	V_{\sigma(x)}Y_\Pi.
	\label{eq:pullback_vertical_fibers}
\end{equation}
Thus
\begin{equation}
	\nabla^\Gamma\sigma
	\in
	\Gamma\bigl(X,\,T^*X\otimes \sigma^*VY_\Pi\bigr).
	\label{eq:covariant_section}
\end{equation}

For the associated bundle
\[
Y_\Pi = Q\times_P T_\Pi,
\]
the vertical bundle has the form
\[
VY_\Pi \cong (Q\times TT_\Pi)/P.
\]
Since \(T_\Pi=\mathbb R^d/\Pi\) is a torus, its tangent bundle is canonically trivial:
\[
TT_\Pi \cong T_\Pi\times \mathbb R^d.
\]
The differential of the \(P\)-action on \(T_\Pi\) acts on the factor \(\mathbb R^d\) by the linear part \(A_p\). 
Hence, after pulling back along \(\sigma\), one obtains the natural isomorphism
\begin{equation}
	\sigma^*VY_\Pi \cong Q\times_P \mathbb R^d.
	\label{eq:sigma_pullback_vertical_trivialized}
\end{equation}
It follows that the covariant differential can be regarded as a section
\begin{equation}
	\nabla^\Gamma\sigma
	\in
	\Gamma\bigl(X,\,T^*X\otimes (Q\times_P \mathbb R^d)\bigr).
	\label{eq:covariant_target_bundle}
\end{equation}
Thus \(\nabla^\Gamma\sigma\) is globally a one-form with values in the associated vector bundle \(Q\times_P\mathbb R^d\); in local trivializations it is represented by an \(\mathbb R^d\)-valued one-form.

\begin{proposition}[Global translational order parameter and covariant differential]
	\label{prop:global_phonon_covariant_differential}
	Let \(Q\to X\) be a fixed principal \(P\)-bundle with \(P\) discrete, and let \(P\) act on \(T_\Pi\) linearly in the symmorphic case or affinely as in \eqref{eq:affine_torus_action} in the nonsymmorphic case. 
	Then the translational order-parameter field on this background is a section
	\[
	\sigma\in\Gamma(X,Y_\Pi),
	\qquad
	Y_\Pi=Q\times_P T_\Pi.
	\]
	Relative to the fixed bundle \(Q\), the bundle \(Y_\Pi\to X\) carries the canonical flat Ehresmann connection \(\Gamma_0\) constructed in Proposition~\ref{prop:ehresmann}. 
	For every section \(\sigma\), its covariant differential is globally defined as
	\begin{equation}
		\nabla^{\Gamma_0}\sigma
		\in
		\Gamma\bigl(X,T^*X\otimes(Q\times_P\mathbb R^d)\bigr).
		\label{eq:main_global_covariant_differential}
	\end{equation}
	In any admissible local trivialization it is represented by the ordinary differential of a local lift, and on overlaps these representatives transform by the linear point-group action.
\end{proposition}

\begin{proof}
	The first statement is the definition of the associated bundle \eqref{eq:Ypi}. 
	The existence of \(\Gamma_0\) follows from Proposition~\ref{prop:ehresmann}. 
	In a trivialization compatible with \(Q\), formula \eqref{eq:flat_covariant_differential} shows that \(\nabla^{\Gamma_0}\sigma\) is represented by \(d\widetilde u_\alpha\), where \(\widetilde u_\alpha\) is a local lift of the torus-valued representative \(\sigma_\alpha\).
	
	On an overlap, the local lifts satisfy
	\[
	\widetilde u_\beta
	=
	A_{g_{\alpha\beta}}^{-1}\widetilde u_\alpha
	+
	a(g_{\alpha\beta}^{-1})
	+
	\lambda_{\alpha\beta},
	\qquad
	\lambda_{\alpha\beta}\in\Pi.
	\]
	The last two terms are locally constant. Hence
	\[
	d\widetilde u_\beta
	=
	A_{g_{\alpha\beta}}^{-1}d\widetilde u_\alpha.
	\]
	These are exactly the transition functions of the associated vector bundle \(Q\times_P\mathbb R^d\). 
	Therefore the local forms \(d\widetilde u_\alpha\) glue to a global section of
	\[
	T^*X\otimes(Q\times_P\mathbb R^d),
	\]
	which proves \eqref{eq:main_global_covariant_differential}.
\end{proof}

\subsection{A minimal twisted example}
\label{subsec:twisted_example}

The following elementary example illustrates why the global bundle language is not merely a change of notation. 
Let \(X=S^1\), let \(\Pi=\mathbb Z\), and let \(T_\Pi=\mathbb R/\mathbb Z\). 
Take \(P=\mathbb Z_2=\{1,r\}\), acting on \(T_\Pi\) by
\[
\rho(r)[u]=[-u].
\]
Let \(Q\to S^1\) be the nontrivial principal \(\mathbb Z_2\)-bundle, equivalently the double cover of the circle. 
Then
\[
Y_\Pi=Q\times_{\mathbb Z_2}(\mathbb R/\mathbb Z)\to S^1
\]
is a twisted torus bundle. 
On two local trivializations whose transition function is \(r\), local representatives of a section satisfy
\[
u_\beta=-u_\alpha\quad \mathrm{mod}\ \mathbb Z.
\]
After choosing local lifts \(\widetilde u_\alpha,\widetilde u_\beta\) to \(\mathbb R\), this becomes
\[
\widetilde u_\beta=-\widetilde u_\alpha+n_{\alpha\beta},
\qquad n_{\alpha\beta}\in\mathbb Z.
\]
Therefore
\[
d\widetilde u_\beta=-d\widetilde u_\alpha.
\]
The local derivatives do not glue to an ordinary global one-form, but they do glue to a section of
\[
T^*S^1\otimes L,
\]
where \(L=Q\times_{\mathbb Z_2}\mathbb R\) is the associated M\"obius line bundle. 
Thus even in this minimal model there need not be a global real-valued displacement field, while the covariant differential \(\nabla^{\Gamma_0}\sigma\) is globally well defined. 
A covariantly constant background section exists at the fixed points \([0]\) and \([1/2]\) of the involution; the corresponding linear phonon field is then a section of \(L\), not a function on \(S^1\).

The jet formulation expresses the bundle geometry in terms of first derivatives. 
The canonical flat connection \(\Gamma_0\) provides the background covariant differential, and \(\nabla^{\Gamma_0}\sigma\) is the global counterpart of the local displacement gradient. 
This is the geometric quantity that enters the first-order Lagrangian for the translational sector.

\subsection{Torus-Valued Field, Local Lift, and Linear Fluctuation}
\label{subsec:toric_lift_linearization}

Before writing the Lagrangian, it is useful to distinguish three objects that are often denoted by the same symbol \(u\) in local physics notation.

First, the global translational order parameter is a section
\[
\sigma\in \Gamma(X,Y_\Pi),
\qquad
Y_\Pi=Q\times_P T_\Pi.
\]
In general, \(\sigma\) is not a map \(X\to\mathbb R^d\), nor is it canonically a map into a fixed torus \(T_\Pi\) without choosing a local trivialization. 
It is this global section to which jets and covariant differentials are applied.
Second, after choosing a local section \(z_\alpha:U_\alpha\to Q\), the global section \(\sigma\) is represented by a torus-valued map
\[
\sigma_\alpha:U_\alpha\to T_\Pi.
\]
On overlaps these representatives satisfy
\[
\sigma_\beta
=
\rho(g_{\alpha\beta}^{-1})\sigma_\alpha.
\]
If local lifts
\[
\widetilde u_\alpha:U_\alpha\to\mathbb R^d,
\qquad
[\widetilde u_\alpha]=\sigma_\alpha,
\]
are chosen, then
\[
\widetilde u_\beta
=
A_{g_{\alpha\beta}}^{-1}\widetilde u_\alpha
+
a(g_{\alpha\beta}^{-1})
+
\lambda_{\alpha\beta},
\qquad
\lambda_{\alpha\beta}\in\Pi.
\]
On each connected component of the overlap, the last two terms are locally constant. Hence the first derivatives of the local lifts transform without the affine shift:
\begin{equation}
	d\widetilde u_\beta
	=
	A_{g_{\alpha\beta}}^{-1}d\widetilde u_\alpha.
	\label{eq:local_lift_derivative_transform}
\end{equation}
This explains why \(\nabla^{\Gamma_0}\sigma\) is globally defined as a section of
\[
T^*X\otimes(Q\times_P\mathbb R^d),
\]
even when the local lifts \(\widetilde u_\alpha\) do not glue to a single \(\mathbb R^d\)-valued field on \(X\).

Third, consider a small fluctuation around a chosen equilibrium section. 
Let \(\sigma_0:X\to Y_\Pi\) be a background section. 
The linear phonon field is then not a new torus-valued order parameter, but a section of the vertical bundle along the background:
\begin{equation}
	\varphi\in\Gamma(X,\sigma_0^*VY_\Pi)
	\simeq
	\Gamma(X,E_\Pi),
	\qquad
	E_\Pi:=Q\times_P\mathbb R^d.
	\label{eq:linear_fluctuation_section}
\end{equation}
In a local trivialization, if \(\widetilde u_{0,\alpha}\) is a local lift of the background, a small configuration may be written as
\begin{equation}
	\widetilde u_\alpha
	=
	\widetilde u_{0,\alpha}+\varphi_\alpha
	\quad
	\mathrm{mod}\ \Pi,
	\label{eq:linearization_torus_exp}
\end{equation}
where the local functions \(\varphi_\alpha:U_\alpha\to\mathbb R^d\) glue linearly:
\[
\varphi_\beta
=
A_{g_{\alpha\beta}}^{-1}\varphi_\alpha.
\]
The affine shifts enter the gluing of the torus-valued representatives themselves, but not the gluing of tangent fluctuations. 
Thus \(\varphi\) is the linear displacement field; in local trivializations it is represented by the ordinary vector-valued displacement field of linear elasticity.

If the background \(\sigma_0\) is covariantly constant with respect to \(\Gamma_0\), then locally
\[
\nabla^{\Gamma_0}\sigma
=
d\varphi_\alpha .
\]
By \eqref{eq:local_lift_derivative_transform}, these local expressions glue to a global section
\begin{equation}
	\nabla^{\Gamma_0}\sigma
	\in
	\Gamma(X,T^*X\otimes E_\Pi).
	\label{eq:flat_cov_diff_lift}
\end{equation}
It is from this linear object that the small-strain tensor, the quadratic elastic energy, and the acoustic phonon spectrum are constructed. 
This agrees with the effective description of phonons as Goldstone modes of broken translations \cite{LeutwylerPhonons,WatanabeMurayama2013,BraunerBook,HayataHidaka2014}, while the absence of independent rotational phonons corresponds to the inverse Higgs mechanism \cite{IvanovOgievetsky1975,LowManohar,BraunerInverseHiggs}.

\section{First-Order Lagrangian Formulation}
\label{sec:first_order_lagrangian}

We now pass from the description of the configuration bundle to the Lagrangian formulation. 
Configurations are sections of
\[
\pi_Y:Y_\Pi\to X.
\]
Therefore a Lagrangian depending on first derivatives is naturally defined on the first jet bundle
\[
J^1Y_\Pi\to X.
\]
Let \(\dim X=n\). 
A first-order Lagrangian on the bundle \(Y_\Pi\to X\) is a fibered morphism
\begin{equation}
	L:J^1Y_\Pi\to \Lambda^nT^*X,
	\label{eq:first_order_lagrangian}
\end{equation}
where \(\Lambda^nT^*X\to X\) is the bundle of \(n\)-forms, interpreted here as Lagrangian densities. 
In local coordinates \((x^\lambda,\tau^i,\tau^i_\lambda)\) on \(J^1Y_\Pi\), induced by a local trivialization of \(Y_\Pi\to X\), the Lagrangian has the form
\begin{equation}
	L=\mathcal L(x^\lambda,\tau^i,\tau^i_\lambda)\,\omega,
	\label{eq:first_order_lagrangian_density}
\end{equation}
where \(\omega=dx^1\wedge\cdots\wedge dx^n\) is the local coordinate volume form on \(X\). 
This is the standard first-order Lagrangian formalism on fiber bundles \cite{SardanashvilyBook,GiachettaBook}.

If \(\sigma:X\to Y_\Pi\) is a section of the configuration bundle, its first jet
\[
J^1\sigma:X\to J^1Y_\Pi
\]
pulls the Lagrangian back to an \(n\)-form \((J^1\sigma)^*L\) on \(X\). 
In local coordinates, if \(\sigma\) is represented by functions \(\tau^i=\sigma^i(x)\), then
\[
(J^1\sigma)^*L
=
\mathcal L\bigl(x^\lambda,\sigma^i(x),\partial_\lambda\sigma^i(x)\bigr)\,\omega.
\]
For a compact domain \(D\subset X\), the corresponding action functional is
\begin{equation}
	S_D[\sigma]=\int_D (J^1\sigma)^*L.
	\label{eq:action_functional}
\end{equation}
Thus a first-order Lagrangian assigns an action to a configuration through its first jet.

\medskip

The Lagrangian \eqref{eq:first_order_lagrangian} has an associated Euler--Lagrange operator
\begin{equation}
	\mathcal E_L:J^2Y_\Pi\to V^*Y_\Pi\otimes \Lambda^nT^*X,
	\label{eq:EL_operator}
\end{equation}
which is locally written as
\begin{equation}
	\mathcal E_L
	=
	\left(
	\frac{\partial \mathcal L}{\partial \tau^i}
	-
	d_\lambda\frac{\partial \mathcal L}{\partial \tau^i_\lambda}
	\right)
	\theta^i\otimes \omega.
	\label{eq:EL_operator_local}
\end{equation}
Here \(d_\lambda\) denotes the total derivative, and
\[
\theta^i=d\tau^i-\tau^i_\lambda dx^\lambda
\]
are the contact \(1\)-forms. The corresponding Euler--Lagrange equations are
\begin{equation}
	\frac{\partial \mathcal L}{\partial \tau^i}
	-
	d_\lambda\frac{\partial \mathcal L}{\partial \tau^i_\lambda}
	=0.
	\label{eq:EL_equations}
\end{equation}
Thus a first-order Lagrangian on \(J^1Y_\Pi\) gives second-order field equations for sections of \(Y_\Pi\to X\).

\medskip

In general, the Lagrangian \eqref{eq:first_order_lagrangian} is an arbitrary fibered morphism on \(J^1Y_\Pi\). 
In the present construction, however, \(Y_\Pi\to X\) carries a connection
\[
\Gamma:Y_\Pi\to J^1Y_\Pi,
\]
and hence a covariant differential
\[
D_\Gamma:J^1Y_\Pi\to \pi_Y^*T^*X\otimes VY_\Pi,
\]
see \eqref{eq:DGamma}. 
It is therefore natural to consider first-order Lagrangians that depend on the first jet only through \(D_\Gamma\). 
Equivalently, we take Lagrangians of the form
\begin{equation}
	L=\widetilde L\circ D_\Gamma,
	\label{eq:lagrangian_factorization}
\end{equation}
where
\begin{equation}
	\widetilde L:
	\pi_Y^*T^*X\otimes VY_\Pi \longrightarrow \Lambda^nT^*X
	\label{eq:Ltilde_map}
\end{equation}
is a fibered morphism over \(X\). 
For the general jet-bundle formalism, see \cite{SaundersJetBundles,GiachettaBook}.

If \(\sigma:X\to Y_\Pi\) is a section, then the factorization \eqref{eq:lagrangian_factorization} means that the Lagrangian depends on the first jet \(J^1\sigma\) only through the covariant differential
\[
\nabla^\Gamma\sigma
=
D_\Gamma\circ J^1\sigma
\in
\Gamma\!\left(X,\,T^*X\otimes \sigma^*VY_\Pi\right),
\]
see \eqref{eq:covariant}. 
In local coordinates this is the dependence on the combinations
\[
\xi^i_\lambda
:=
\tau^i_\lambda-\Gamma^i_\lambda(x,\tau).
\]
Accordingly, the local Lagrangian density has the form
\begin{equation}
	L=
	\widetilde{\mathcal L}(x^\lambda,\tau^i,\xi^i_\lambda)\,\omega,
	\qquad
	\xi^i_\lambda=\tau^i_\lambda-\Gamma^i_\lambda(x,\tau).
	\label{eq:factorized_lagrangian_density}
\end{equation}

For the translational sector, this is the natural class of first-order Lagrangians: \(\nabla^\Gamma\sigma\) is the global replacement for the local displacement gradient.

\medskip

For a Lagrangian of the form \eqref{eq:factorized_lagrangian_density}, the Euler--Lagrange equations \eqref{eq:EL_equations} become
\begin{equation}
	\frac{\partial \widetilde{\mathcal L}}{\partial \tau^i}
	-
	\frac{\partial \widetilde{\mathcal L}}{\partial \xi^j_\lambda}
	\frac{\partial \Gamma^j_\lambda}{\partial \tau^i}
	-
	d_\lambda\!\left(
	\frac{\partial \widetilde{\mathcal L}}{\partial \xi^i_\lambda}
	\right)
	=0.
	\label{eq:EL_factorized}
\end{equation}
For the canonical flat connection \(\Gamma_0\) constructed in Sec.~\ref{subsec:ehresmann}, one has \(\Gamma^i_{0\,\lambda}=0\) in admissible local trivializations, and therefore
\[
\xi^i_\lambda=\tau^i_\lambda.
\]
For \(\Gamma_0\), the Lagrangian depends locally on the ordinary first derivatives of the field. 
After imposing the usual elastic symmetry conditions, in particular dependence of the spatial energy on the symmetrized displacement gradient, quadratic Lagrangians reproduce the standard local harmonic models of linear elasticity and acoustic phonons. 
The role of the present construction is to place these local equations on the global configuration bundle \(Y_\Pi\to X\).

\subsection{Quadratic Lagrangians for the Translational Sector}
\label{subsec:quadratic_phonon_lagrangians}

At low energy, the natural first approximation is a Lagrangian quadratic in the covariant differential. 
The Euler--Lagrange operator for a general first-order Lagrangian was introduced above, see \eqref{eq:EL_operator_local}--\eqref{eq:EL_equations}. 
We now specialize it to Lagrangians that factor through
\[
D_\Gamma:J^1Y_\Pi\to \pi_Y^*T^*X\otimes VY_\Pi.
\]

Let
\begin{equation}
	\mathbf G:
	\left(\pi_Y^*T^*X\otimes VY_\Pi\right)
	\times_{Y_\Pi}
	\left(\pi_Y^*T^*X\otimes VY_\Pi\right)
	\longrightarrow \mathbb R
	\label{eq:bundle_metric_quadratic}
\end{equation}
be a smooth symmetric bilinear form on the vector bundle
\[
\pi_Y^*T^*X\otimes VY_\Pi \to Y_\Pi.
\]
After choosing a volume form \(\omega\) on \(X\), this bilinear form defines the quadratic Lagrangian
\begin{equation}
	L_{\mathbf G}
	=
	\frac12\,\mathbf G(D_\Gamma,D_\Gamma)\,\omega.
	\label{eq:quadratic_lagrangian_global}
\end{equation}
Here \(\mathbf G(D_\Gamma,D_\Gamma)\) denotes the function on \(J^1Y_\Pi\) obtained by applying \(\mathbf G\) to the covariant differential \(D_\Gamma\).

If \(\sigma:X\to Y_\Pi\) is a section, then the induced Lagrangian density is
\begin{equation}
	(J^1\sigma)^*L_{\mathbf G}
	=
	\frac12\,(\sigma^*\mathbf G)
	\bigl(\nabla^\Gamma\sigma,\nabla^\Gamma\sigma\bigr)\,\omega,
	\label{eq:quadratic_lagrangian_on_section}
\end{equation}
where
\[
\nabla^\Gamma\sigma
=
D_\Gamma\circ J^1\sigma
\in
\Gamma\!\left(X,\,T^*X\otimes \sigma^*VY_\Pi\right).
\]
Thus the Lagrangian depends on the section \(\sigma\) through its first covariant differential.

In local coordinates \((x^\lambda,\tau^i,\tau^i_\lambda)\) on \(J^1Y_\Pi\), set
\begin{equation}
	\xi^i_\lambda:=\tau^i_\lambda-\Gamma^i_\lambda(x,\tau).
	\label{eq:xi_definition_quadratic}
\end{equation}
Then the Lagrangian \eqref{eq:quadratic_lagrangian_global} is written as
\begin{equation}
	L_{\mathbf G}
	=
	\frac12\,G^{\lambda\mu}_{ij}(x,\tau)\,
	\xi^i_\lambda \xi^j_\mu\,\omega,
	\label{eq:quadratic_lagrangian_local}
\end{equation}
where the coefficients \(G^{\lambda\mu}_{ij}\) are symmetric under the simultaneous interchange
\[
(\lambda,i)\longleftrightarrow(\mu,j).
\]
These local coefficients are not arbitrary functions on \(\mathbb R^d\). In local lifts of the torus coordinate they are periodic with respect to the lattice,
\[
G^{\lambda\mu}_{ij}(x,\tau+\ell)=G^{\lambda\mu}_{ij}(x,\tau),
\qquad \ell\in\Pi,
\]
and on overlaps they satisfy the tensorial \(P\)-equivariance rule induced by the affine action on \(T_\Pi\) and its linear differential on vertical tangent vectors. Explicitly, if
\[
\tau_\beta=\rho(g_{\alpha\beta}^{-1})\tau_\alpha,
\qquad
\xi_\beta=A_{g_{\alpha\beta}}^{-1}\xi_\alpha,
\]
then, writing \(B=A_{g_{\alpha\beta}}^{-1}\), one has
\[
G^{\lambda\mu}_{\alpha,ij}(x,\tau_\alpha)
=
B^r{}_{i}B^s{}_{j}
G^{\lambda\mu}_{\beta,rs}(x,\tau_\beta).
\]
These periodicity and equivariance conditions are precisely what make \(\mathbf G\) a globally defined bilinear form on \(\pi_Y^*T^*X\otimes VY_\Pi\).

If \(\sigma\) is represented locally by functions \(\tau^i=\sigma^i(x)\), then
\begin{equation}
	(J^1\sigma)^*L_{\mathbf G}
	=
	\frac12\,G^{\lambda\mu}_{ij}(x,\sigma(x))\,
	\nabla^\Gamma_\lambda\sigma^i\,
	\nabla^\Gamma_\mu\sigma^j\,\omega,
	\label{eq:quadratic_lagrangian_section_local}
\end{equation}
where
\[
\nabla^\Gamma_\lambda\sigma^i
=
\partial_\lambda\sigma^i-\Gamma^i_\lambda(x,\sigma(x)).
\]

Lagrangians of the form \eqref{eq:quadratic_lagrangian_global} are the bundle-theoretic analogue of the harmonic approximation in elasticity: they describe low-energy fluctuations of translational order through the first covariant differential. 
The use of \(D_\Gamma\), rather than ordinary partial derivatives, is what makes the formulation global on
\[
Y_\Pi=Q\times_P T_\Pi.
\]

At the same time, a general symmetric bilinear form \(\mathbf G\) on \(\pi_Y^*T^*X\otimes VY_\Pi\) is more general than an ordinary elastic energy. 
The Goldstone, or acoustic, interpretation requires a derivative-only leading theory: in adapted local trivializations the leading Lagrangian is invariant under constant translations of the torus coordinate and contains no potential term for \(\tau\). 
Equivalently, the coefficients relevant for the homogeneous harmonic approximation are independent of the absolute torus phase, up to the prescribed \(P\)-equivariance. 
If explicit \(\tau\)-dependent terms are retained, they should be interpreted as pinning or other controlled symmetry-breaking effects, and the gapless acoustic conclusion need not follow. 
To recover the acoustic phonon model after eliminating independent rotational Goldstone variables, one also imposes the usual objectivity condition: in the leading spatial sector the energy depends on the symmetrized displacement gradient, namely on the small-strain tensor, and not on a pure infinitesimal rigid rotation. 
This condition is implemented explicitly in Appendix~\ref{app:cubic_example} through \(\varepsilon_{ij}\) and the standard elastic constants.

\medskip

We now derive the equations of motion for the quadratic Lagrangian. 
Using the notation
\[
\xi^i_\lambda=\tau^i_\lambda-\Gamma^i_\lambda(x,\tau),
\]
write
\begin{equation}
	\mathcal L_{\mathbf G}
	=
	\frac12\,G^{\lambda\mu}_{ij}(x,\tau)\,
	\xi^i_\lambda \xi^j_\mu .
	\label{eq:quadratic_density_xi}
\end{equation}
Assume that \(G^{\lambda\mu}_{ij}\) is symmetric under the simultaneous interchange
\[
(\lambda,i)\longleftrightarrow(\mu,j).
\]
Then
\begin{equation}
	\frac{\partial \mathcal L_{\mathbf G}}{\partial \tau^k_\nu}
	=
	G^{\nu\mu}_{kj}(x,\tau)\,\xi^j_\mu,
	\label{eq:quadratic_first_momentum}
\end{equation}
and
\begin{equation}
	\frac{\partial \mathcal L_{\mathbf G}}{\partial \tau^k}
	=
	\frac12\,
	\frac{\partial G^{\lambda\mu}_{ij}}{\partial \tau^k}\,
	\xi^i_\lambda \xi^j_\mu
	-
	G^{\lambda\mu}_{ij}\,
	\frac{\partial \Gamma^i_\lambda}{\partial \tau^k}\,
	\xi^j_\mu .
	\label{eq:quadratic_first_tau}
\end{equation}
Substitution into the Euler--Lagrange equations gives
\begin{equation}
	\frac12\,
	\frac{\partial G^{\lambda\mu}_{ij}}{\partial \tau^k}\,
	\xi^i_\lambda \xi^j_\mu
	-
	G^{\lambda\mu}_{ij}\,
	\frac{\partial \Gamma^i_\lambda}{\partial \tau^k}\,
	\xi^j_\mu
	-
	d_\nu\!\left(
	G^{\nu\mu}_{kj}\,\xi^j_\mu
	\right)
	=0.
	\label{eq:quadratic_EL_explicit}
\end{equation}

The structure of \eqref{eq:quadratic_EL_explicit} shows that the leading second-order term is
\[
-\,G^{\nu\mu}_{kj}(x,\tau)\,\tau^j_{\nu\mu}.
\]
All remaining terms contain at most first derivatives of the field, together with derivatives of the coefficients \(G^{\lambda\mu}_{ij}\) and \(\Gamma^i_\lambda\). 
Thus a quadratic Lagrangian of the above form gives, in general, a quasilinear second-order system. 
Its principal symbol determines the type of the system and the propagation of small disturbances.
The system takes a particularly simple form in an admissible local trivialization for the canonical flat Ehresmann connection \(\Gamma_0\), where
\[
\Gamma^i_{0\,\lambda}=0.
\]
Then
\[
\xi^i_\lambda=\tau^i_\lambda,
\]
and the quadratic Lagrangian becomes
\begin{equation}
	L_{\mathbf G}
	=
	\frac12\,G^{\lambda\mu}_{ij}(x,\tau)\,
	\tau^i_\lambda\tau^j_\mu\,\omega.
	\label{eq:quadratic_lagrangian_flat_trivialization}
\end{equation}
The Euler--Lagrange equations reduce to
\begin{equation}
	\frac12\,
	\frac{\partial G^{\lambda\mu}_{ij}}{\partial \tau^k}\,
	\tau^i_\lambda\tau^j_\mu
	-
	d_\nu\!\left(
	G^{\nu\mu}_{kj}\,\tau^j_\mu
	\right)
	=0.
	\label{eq:quadratic_EL_flat_explicit}
\end{equation}
If the coefficients \(G^{\lambda\mu}_{ij}\) are independent of \(\tau\), this becomes
\begin{equation}
	d_\nu\!\left(
	G^{\nu\mu}_{kj}(x)\,\tau^j_\mu
	\right)=0.
	\label{eq:quadratic_EL_tau_independent}
\end{equation}
For constant coefficients one obtains the linear second-order system
\begin{equation}
	G^{\nu\mu}_{kj}\,\tau^j_{\nu\mu}=0.
	\label{eq:quadratic_EL_constant_coefficients}
\end{equation}

The leading part of the Euler--Lagrange equations corresponding to 
\eqref{eq:quadratic_lagrangian_global} is
\[
-\,G^{\nu\mu}_{kj}(x,\tau)\,\tau^j_{\nu\mu}.
\]
Thus, at a covector \(\zeta\in T_x^*X\), the principal symbol is the matrix
\begin{equation}
	\sigma_{\mathrm{pr}}(\zeta)^k{}_j
	=
	-\,G^{\nu\mu}_{kj}(x,\tau)\,\zeta_\nu\zeta_\mu.
	\label{eq:principal_symbol_quadratic}
\end{equation}
Up to this conventional overall sign, the principal symbol is determined by the bilinear form \(\mathbf G\). 
This symbol controls the type of the linearized system and the propagation of small disturbances.

\subsection{Noether currents in brief}
\label{subsec:noether_brief}

The same first-order formalism also identifies the conserved currents associated with variational symmetries. 
For a Lagrangian density \(L=\mathcal L\omega\), set
\begin{equation}
	\pi_i^\lambda:=\frac{\partial\mathcal L}{\partial \tau^i_\lambda}.
	\label{eq:canonical_momenta_brief}
\end{equation}
For the quadratic Lagrangian \eqref{eq:quadratic_lagrangian_local},
\begin{equation}
	\pi_i^\lambda
	=
	G^{\lambda\mu}_{ij}(x,\tau)\,\xi^j_\mu,
	\qquad
	\xi^j_\mu=\tau^j_\mu-\Gamma^j_\mu(x,\tau).
	\label{eq:quadratic_momenta_brief}
\end{equation}
Thus the momenta are built from the covariant differential of the translational order parameter. 
When the Lagrangian is invariant under the corresponding infinitesimal torus shift, for example in the homogeneous derivative-only case with \(a \in \Gamma(X,Q\times_P\mathbb R^d)\) represented locally by a constant or covariantly constant shift, the corresponding internal current is locally
\begin{equation}
	J^\lambda_{(a)}=\pi_i^\lambda a^i
	\label{eq:internal_current_brief}
\end{equation}
up to the usual improvement or boundary term. 
Similarly, if a vector field \(\zeta=\zeta^\mu\partial_\mu\) on the base has a horizontal lift that is a variational symmetry, the associated canonical energy-momentum tensor is
\begin{equation}
	T^\lambda{}_{\mu}
	=
	\pi_i^\lambda\bigl(\tau^i_\mu-\Gamma^i_\mu\bigr)
	-
	\delta^\lambda_\mu\mathcal L.
	\label{eq:canonical_energy_momentum_tensor}
\end{equation}
On solutions, the corresponding currents satisfy the standard weak conservation laws. 
These formulae are local representatives of globally defined associated-bundle objects; under a change of crystallographic trivialization they transform by the same \(P\)-representation as the covariant differential.

\subsection{Linear Approximation and Dispersion Relations}
\label{subsec:linear_dispersion}

To obtain the local phonon spectrum, consider small perturbations around an equilibrium section. Let
\begin{equation}
	\sigma_0:X\to Y_\Pi
	\label{eq:equilibrium_section}
\end{equation}
be a smooth solution of the equations of motion satisfying
\begin{equation}
	\nabla^{\Gamma_0}\sigma_0=0,
	\label{eq:equilibrium_covariantly_constant}
\end{equation}
where \(\Gamma_0\) is the canonical flat connection constructed above. 
Geometrically, \(\sigma_0\) represents an undeformed reference configuration on the fixed crystallographic background. 
The existence of such a global covariantly constant section is not automatic: it requires the holonomy action of the flat torus bundle to fix a point of \(T_\Pi\). 
If this condition fails, the linearization below should be understood locally, for example on a simply connected patch where a background lift has been chosen.

In a local trivialization, the translational order parameter is represented by a map into the torus
\[
T_\Pi=\mathbb R^d/\Pi.
\]
Therefore linearization is a local operation. 
Choose a local chart on \(T_\Pi\) near the image of the background representative \(\sigma_{0,\alpha}\), and write
\begin{equation}
	\tau^i=\tau_0^i+\varphi^i,
	\label{eq:small_fluctuation}
\end{equation}
where \(\varphi=(\varphi^i)\) is a small \(\mathbb R^d\)-valued fluctuation in this chart. 
In a local trivialization compatible with the canonical flat Ehresmann connection,
\[
\Gamma^i_{0\,\lambda}=0.
\]
Hence \(\nabla^{\Gamma_0}\) is represented locally by the ordinary first derivative with respect to the base coordinates. 
In particular, the condition \eqref{eq:equilibrium_covariantly_constant} implies that the local coordinates \(\tau_0^i\) of the equilibrium section are constant.

\medskip

Consider the quadratic Lagrangian
\begin{equation}
	L_{\mathbf G}
	=
	\frac12\,G^{\lambda\mu}_{ij}(x,\tau)\,
	\tau^i_\lambda\tau^j_\mu\,\omega,
	\label{eq:quadratic_lagrangian_for_linearization}
\end{equation}
written in an admissible local trivialization for the canonical flat connection \(\Gamma_0\). 
Expanding around the covariantly constant background \(\sigma_0\), with
\[
\tau^i=\tau_0^i+\varphi^i,
\]
gives the quadratic part
\begin{equation}
	L^{(2)}
	=
	\frac12\,G^{\lambda\mu}_{ij}(x,\tau_0)\,
	\partial_\lambda\varphi^i\,\partial_\mu\varphi^j\,\omega.
	\label{eq:linearized_lagrangian_general}
\end{equation}
This quadratic Lagrangian determines the linearized equations of motion.

If the equilibrium configuration is locally homogeneous, the coefficients 
\(G^{\lambda\mu}_{ij}(x,\tau_0)\) may be replaced by their constant values. 
Then
\begin{equation}
	L^{(2)}
	=
	\frac12\,G^{\lambda\mu}_{ij}\,
	\partial_\lambda\varphi^i\,\partial_\mu\varphi^j\,\omega,
	\label{eq:linearized_lagrangian_constant}
\end{equation}
where \(G^{\lambda\mu}_{ij}\) are constant coefficients symmetric under the simultaneous interchange
\[
(\lambda,i)\longleftrightarrow (\mu,j).
\]

The corresponding linearized Euler--Lagrange equations are
\begin{equation}
	G^{\lambda\mu}_{ij}\,\partial_\lambda\partial_\mu\varphi^j=0.
	\label{eq:linearized_EL_general}
\end{equation}
Thus one obtains a linear second-order system for the small displacement fluctuations. 
At a covector \(\zeta=\zeta_\lambda dx^\lambda\in T_x^*X\), its principal symbol is the matrix
\begin{equation}
	\sigma_{\mathrm{pr}}(\zeta)^i{}_j
	=
	G^{\lambda\mu}_{ij}\,\zeta_\lambda\zeta_\mu.
	\label{eq:principal_symbol_linearized}
\end{equation}
This symbol determines the type of the linearized system and its dispersion properties.

\medskip

In the dynamical setting
\[
X=\mathbb R_t\times M,
\]
write local coordinates as \(x^\lambda=(t,x^a)\), where \(a=1,\dots,d\). 
Then the density corresponding to \eqref{eq:linearized_lagrangian_constant} can be decomposed as
\begin{equation}
	\mathcal L^{(2)}
	=
	\frac12\,\rho_{ij}\,\partial_t\varphi^i\,\partial_t\varphi^j
	+
	K^a_{ij}\,\partial_t\varphi^i\,\partial_a\varphi^j
	-
	\frac12\,C^{ab}_{ij}\,\partial_a\varphi^i\,\partial_b\varphi^j,
	\label{eq:linearized_lagrangian_split}
\end{equation}
where
\[
\rho_{ij}=G^{tt}_{ij},
\qquad
K^a_{ij}=G^{ta}_{ij},
\qquad
C^{ab}_{ij}=-\,G^{ab}_{ij}.
\]
The matrix \(\rho_{ij}\) is the mass-density matrix, while \(C^{ab}_{ij}\) gives the spatial elastic coefficients in the linear approximation.

If, in addition, the theory is invariant under time reversal, with \(\varphi\) even and \(\partial_t\varphi\) odd, then the mixed terms linear in \(\partial_t\varphi\) are excluded. 
Thus \(K^a_{ij}=0\), and the quadratic Lagrangian density reduces to
\begin{equation}
	\mathcal L^{(2)}
	=
	\frac12\,\rho_{ij}\,\partial_t\varphi^i\,\partial_t\varphi^j
	-
	\frac12\,C^{ab}_{ij}\,\partial_a\varphi^i\,\partial_b\varphi^j.
	\label{eq:linearized_lagrangian_time_reversal}
\end{equation}
The corresponding equations of motion are
\begin{equation}
	\rho_{ij}\,\partial_t^2\varphi^j
	-
	C^{ab}_{ij}\,\partial_a\partial_b\varphi^j
	=
	0.
	\label{eq:linearized_phonon_equation}
\end{equation}

\medskip

To obtain the local dispersion relations, look for plane-wave solutions of the form
\begin{equation}
	\varphi^i(t,x)=\alpha^i\exp\!\bigl(\mathrm i(k_ax^a-\omega t)\bigr),
	\label{eq:plane_wave_ansatz}
\end{equation}
where \(\alpha^i\) is a constant polarization amplitude.
Substitution of \eqref{eq:plane_wave_ansatz} into \eqref{eq:linearized_phonon_equation} gives the algebraic system
\begin{equation}
	\bigl(\mathcal C_{ij}(k)-\omega^2\rho_{ij}\bigr)\alpha^j=0,
	\label{eq:christoffel_equation}
\end{equation}
where
\begin{equation}
	\mathcal C_{ij}(k):=C^{ab}_{ij}\,k_ak_b
	\label{eq:christoffel_matrix}
\end{equation}
is the Christoffel, or acoustic, matrix associated with the wave vector \(k\) \cite{NyePhysicalProperties,MusgraveCrystalAcoustics}.
The notation \(\mathcal C\) is used to avoid confusion with the connection \(\Gamma_0\).
Nontrivial solutions exist if and only if
\begin{equation}
	\det\bigl(\mathcal C(k)-\omega^2\rho\bigr)=0.
	\label{eq:dispersion_relation_general}
\end{equation}
This is the local dispersion equation for the acoustic phonon modes. 
Equivalently, one has the generalized eigenvalue problem
\begin{equation}
	\mathcal C(k)\,\alpha=\omega^2\rho\,\alpha.
	\label{eq:eigenvalue_problem_phonons}
\end{equation}
If \(\rho\) is nondegenerate, this may be written as
\[
\rho^{-1}\mathcal C(k)\,\alpha=\omega^2\alpha.
\]
Since
\[
\mathcal C(k)=|k|^2\mathcal C(\hat k),
\qquad
\hat k=\frac{k}{|k|},
\]
the long-wavelength dispersion is acoustic:
\begin{equation}
	\omega_\alpha(k)=c_\alpha(\hat k)\,|k|.
	\label{eq:acoustic_dispersion}
\end{equation}
Here \(c_\alpha(\hat k)^2\) are the eigenvalues of \(\rho^{-1}\mathcal C(\hat k)\), and \(c_\alpha(\hat k)\) are the corresponding propagation velocities.

\medskip

In the anisotropic case the coefficients of the linearized Lagrangian are not arbitrary. 
They are constrained by the point group \(P\) of the crystal. 
Since the configuration bundle is
\[
Y_\Pi=Q\times_P T_\Pi,
\]
the tensorial quantities entering the Lagrangian must be compatible with the \(P\)-action. 
Equivalently, in local frames adapted to the \(P\)-reduction \(Q\), the coefficients must transform according to the point-group representation. 
In particular, the bilinear form \(\mathbf G\), and hence the induced coefficients \(\rho_{ij}\) and \(C^{ab}_{ij}\), are constrained by the point-group symmetry.
In a local orthonormal frame adapted to the reduction
\[
Q\subset F_g(M),
\]
the point group acts on both spatial indices and displacement indices. 
Thus, for every \(R\in P\subset O(d)\), the coefficients satisfy
\begin{equation}
	\rho_{ij}
	=
	R_i{}^k\,R_j{}^l\,\rho_{kl},
	\label{eq:P_invariance_density}
\end{equation}
and
\begin{equation}
	C^{ab}_{ij}
	=
	R^a{}_c\,R^b{}_d\,R_i{}^k\,R_j{}^l\,C^{cd}_{kl}.
	\label{eq:P_invariance_elastic_tensor}
\end{equation}
Therefore the point group \(P\) selects the admissible components of the elastic coefficients and determines the allowed anisotropy of the dispersion relation \eqref{eq:dispersion_relation_general}.

Equivalently, the Christoffel matrix \eqref{eq:christoffel_matrix} satisfies the covariance condition
\begin{equation}
	\mathcal C(Rk)=R\,\mathcal C(k)\,R^{-1},
	\qquad R\in P.
	\label{eq:christoffel_covariance}
\end{equation}
Together with the \(P\)-invariance of \(\rho\), this implies that the generalized eigenvalues of the pair \((\mathcal C(k),\rho)\), and hence the unordered set of phonon frequencies \(\omega_\alpha(k)\), are constant on point-group orbits in wave-vector space. 
Thus the point group determines the allowed anisotropy of the sound velocities.

\medskip

In the isotropic limit, the coefficients are invariant under the full orthogonal group \(O(d)\). 
After choosing a local orthonormal frame and identifying spatial and internal indices, one may write
\begin{equation}
	\rho_{ij}=\rho\,\delta_{ij},
	\label{eq:isotropic_density}
\end{equation}
and
\begin{equation}
	C^{ab}_{ij}
	=
	\lambda\,\delta^a_i\delta^b_j
	+
	\mu\bigl(\delta_{ij}\delta^{ab}+\delta^a_j\delta^b_i\bigr),
	\label{eq:isotropic_elasticity_tensor}
\end{equation}
where \(\lambda\) and \(\mu\) are the Lam\'e coefficients. 
Then the Christoffel matrix becomes
\begin{equation}
	\mathcal C_{ij}(k)
	=
	\mu |k|^2\delta_{ij}
	+
	(\lambda+\mu)\,k_i k_j.
	\label{eq:isotropic_christoffel}
\end{equation}
The dispersion equation therefore splits into one longitudinal branch and \(d-1\) transverse branches:
\begin{equation}
	\omega_L^2
	=
	\frac{\lambda+2\mu}{\rho}\,|k|^2,
	\qquad
	\omega_T^2
	=
	\frac{\mu}{\rho}\,|k|^2.
	\label{eq:longitudinal_transverse_dispersion}
\end{equation}
The longitudinal mode has polarization parallel to \(k\), while the transverse modes have polarizations orthogonal to \(k\).

\medskip

The quadratic fluctuation theory obtained from the associated bundle \(Y_\Pi\to X\) therefore reproduces the standard local acoustic spectrum. 
The additional information retained by the geometric formulation is global: the translational order parameter is torus-valued, the linearized fields are sections of the associated vector bundle \(Q\times_P\mathbb R^d\), and the coefficients entering the dispersion relation must be compatible with the point-group action. 
Thus the point-group reduction does not change the local form of the acoustic equations, but it constrains the allowed anisotropy and the gluing of phonon modes across different crystallographic trivializations.

\section{Conclusion}
\label{sec:conclusion}

This paper has formulated the translational, or phonon, sector of a crystalline phase as a global field theory on an associated torus bundle. 
The orientational part of the crystalline order was kept fixed, in the form of a \(P\)-reduction \(Q\subset F_g(M)\). 
With this crystallographic background fixed, the translational order parameter is not, in general, a globally defined \(\mathbb R^d\)-valued displacement field. 
It is a section of
\[
Y_\Pi=Q\times_P T_\Pi,
\qquad
T_\Pi=\mathbb R^d/\Pi,
\]
the associated torus bundle constructed above.

The formulation also records the distinction between symmorphic and nonsymmorphic crystallographic groups. 
In the symmorphic case, the point group acts linearly on \(T_\Pi\). 
In the nonsymmorphic case, the action is affine, as in \eqref{eq:affine_torus_action}. 
The affine shifts enter the gluing of local torus-valued representatives of the order parameter, whereas the vertical linearized theory depends only on the differential of this action, namely on the usual linear point-group action.

The central geometric step was the construction of a flat Ehresmann connection on \(Y_\Pi\to X\), canonical relative to the fixed \(P\)-bundle \(Q\). 
It exists because the structure group is discrete: in admissible local trivializations the transition functions are locally constant, so the horizontal spaces \(TU_\alpha\oplus0\) glue globally. 
In jet language this connection is a section of \(J^1Y_\Pi\to Y_\Pi\). 
Proposition~\ref{prop:global_phonon_covariant_differential} gives the corresponding covariant differential,
\[
\nabla^{\Gamma_0}\sigma\in
\Gamma\bigl(X,T^*X\otimes(Q\times_P\mathbb R^d)\bigr).
\]
For the canonical flat connection, it is represented locally by the ordinary derivative of a displacement lift, but globally it is a well-defined section of the associated vector bundle \(T^*X\otimes(Q\times_P\mathbb R^d)\). 
The twisted circle example in Sec.~\ref{subsec:twisted_example} shows explicitly how this can hold even when no ordinary global displacement function exists.

Using this covariant differential, we wrote a first-order Lagrangian formulation on \(J^1Y_\Pi\). 
The most relevant low-energy class consists of Lagrangians that factor through \(D_\Gamma\) and are quadratic in the covariant differential. 
Their Euler--Lagrange equations give a quasilinear second-order system in general and reduce, for constant coefficients in an admissible flat trivialization, to the standard local linear equations of elasticity.

The Noether currents fit the same picture. 
Whenever internal torus shifts or symmetries of the base \(X\) are variational symmetries of the chosen Lagrangian, the corresponding local currents are built from the canonical momenta and, for base symmetries, from the canonical energy-momentum tensor \eqref{eq:canonical_energy_momentum_tensor}. 
They transform according to the same associated-bundle structures as the covariant differential.

The linearized theory recovers the usual local acoustic phonon spectrum. 
The Christoffel matrix, the dispersion relation, and, in the isotropic limit, the longitudinal/transverse decomposition all take their standard forms. 
What the bundle construction adds is the global interpretation of these objects. 
The translational order parameter is torus-valued from the start, the linearized fields are sections of \(Q\times_P\mathbb R^d\), and the coefficients in the quadratic theory must be compatible with the point-group action. 
Thus the point-group reduction does not change the local acoustic equations, but it constrains the allowed anisotropy of the sound velocities and the gluing of phonon modes across crystallographic trivializations.

Appendix~\ref{app:cubic_example} checks that the formalism reproduces the textbook three-constant theory of a cubic crystal.

The main limitation of the paper is deliberate: the orientational sector is fixed, and the translational field is assumed to be smooth in the main text. 
A fuller theory would allow defects such as dislocations and disclinations from the beginning, treat the \(P\)-reduction dynamically, and incorporate additional background structures such as magnetic or time-reversing symmetries. 
The construction presented here should be viewed as the smooth background layer of such a theory. 
It identifies the global configuration space for the translational order parameter and the covariant differential that replaces the local displacement gradient.

\appendix


\section{Local Cubic Elasticity from the Global Translational Field}
\label{app:cubic_example}

This appendix gives a compact local check of the global construction for a three-dimensional cubic crystal. 
Let \(X=\mathbb R_t\times M\), \(\dim M=3\), and let \(Q_M\subset F_g(M)\) be a reduction to the cubic point group \(P=O_h\subset O(3)\), or to its orientation-preserving subgroup. 
In the dynamical setting we use \(Q=\mathrm{pr}_M^*Q_M\to X\). 
After fixing units in which the cubic lattice spacing is one, take \(\Pi\simeq\mathbb Z^3\), \(T_\Pi=\mathbb R^3/\Pi\), and
\[
Y_\Pi=Q\times_P T_\Pi\to X.
\]
Thus the translational order parameter remains the global section \(\sigma:X\to Y_\Pi\).

Linearization around a covariantly constant background section gives the associated vector bundle \(E_\Pi:=Q\times_P\mathbb R^3\to X\). 
The linear displacement field is therefore globally a section \(\varphi\in\Gamma(X,E_\Pi)\), not a globally defined \(\mathbb R^3\)-valued function unless \(E_\Pi\) is trivial. 
Since \(Q_M\subset F_g(M)\), there is a natural identification \(E_\Pi\simeq\mathrm{pr}_M^*TM\), given in a cubic frame by \([q,v]\mapsto q(v)\).

Let
\[
\beta:=\nabla^{\Gamma_0}\sigma
\]
denote the covariant differential of the translational field. 
After linearization around the background section, it is regarded as a section
\[
\beta\in
\Gamma\bigl(X,T^*X\otimes E_\Pi\bigr).
\]
Its temporal part \(v:=\beta(\partial_t)\in\Gamma(X,E_\Pi)\) is the velocity field, while its spatial part
\[
F:=\beta|_{\mathrm{pr}_M^*TM}
\in
\Gamma\bigl(X,\mathrm{pr}_M^*T^*M\otimes E_\Pi\bigr)
\]
is the global displacement-gradient object. 
Using \(E_\Pi\simeq\mathrm{pr}_M^*TM\), define the global small-strain tensor by
\[
\varepsilon:=\operatorname{sym}(F^\flat)
\in
\Gamma\bigl(X,\mathrm{pr}_M^*\mathrm{Sym}^2T^*M\bigr).
\]
In an admissible crystallographic trivialization over a contractible open set, \(\varphi\) is represented by the usual local displacement \(u=(u^1,u^2,u^3)\), and
\[
\beta=\partial_\lambda u^i\,dx^\lambda\otimes e_i,
\qquad
\varepsilon_{ij}=\frac12(\partial_i u_j+\partial_j u_i).
\]

The cubic elastic energy is the local representative of a \(P\)-invariant quadratic form \(W_Q\) on \(\mathrm{pr}_M^*\mathrm{Sym}^2T^*M\). 
In a cubic frame this representative is
\[
W_{\mathrm{cub}}(\varepsilon)
=
\frac12 C_{11}(\varepsilon_{11}^2+\varepsilon_{22}^2+\varepsilon_{33}^2)
+
C_{12}(\varepsilon_{11}\varepsilon_{22}
+\varepsilon_{22}\varepsilon_{33}
+\varepsilon_{33}\varepsilon_{11})
+
2C_{44}(\varepsilon_{12}^2+\varepsilon_{23}^2+\varepsilon_{31}^2).
\]
The corresponding quadratic Lagrangian density is globally
\[
\mathcal L^{(2)}
=
\frac12\rho\,\langle v,v\rangle-W_Q(\varepsilon),
\]
and in the above local frame it becomes the standard cubic elastic Lagrangian. 
The stress tensor is defined by \(\delta W_{\mathrm{cub}}=\sigma_{ij}\delta\varepsilon_{ij}\), hence
\[
\sigma_{ii}=C_{11}\varepsilon_{ii}+C_{12}\sum_{j\neq i}\varepsilon_{jj},
\qquad
\sigma_{ij}=2C_{44}\varepsilon_{ij}\quad(i\neq j),
\]
and the local equations of motion are
\[
\rho\,\partial_t^2u_i=\partial_j\sigma_{ij}.
\]

For a plane wave \(u_i=\alpha_i\exp\!\bigl(\mathrm i(k\cdot x-\omega t)\bigr)\), the equations become
\[
\mathcal C_{ij}(k)\alpha_j=\rho\omega^2\alpha_i,
\]
where the cubic Christoffel matrix \cite{NyePhysicalProperties,MusgraveCrystalAcoustics} is compactly written as
\[
\mathcal C_{ij}(k)
=
C_{44}|k|^2\delta_{ij}
+
(C_{12}+C_{44})k_i k_j
+
(C_{11}-C_{12}-2C_{44})k_i^2\delta_{ij},
\]
with no summation over \(i\) in the last term. 
This gives, for example, in the \([100]\) direction
\[
\omega_L^2=\frac{C_{11}}{\rho}k^2,
\qquad
\omega_T^2=\frac{C_{44}}{\rho}k^2,
\]
with a doubly degenerate transverse branch, while in the \([111]\) direction
\[
\omega_L^2=\frac{C_{11}+2C_{12}+4C_{44}}{3\rho}k^2,
\qquad
\omega_T^2=\frac{C_{11}-C_{12}+C_{44}}{3\rho}k^2.
\]
The standard positivity conditions are \(C_{44}>0\), \(C_{11}-C_{12}>0\), \(C_{11}+2C_{12}>0\), and \(\rho>0\).

Thus the usual cubic acoustic equations are recovered as local component expressions of globally defined objects: the fluctuation field is a section of \(E_\Pi\), the velocity and displacement gradient are the temporal and spatial parts of the local representative of \(\nabla^{\Gamma_0}\sigma\), and the cubic elastic energy is the local representative of a \(P\)-invariant quadratic form determined by the reduction \(Q\).



\input{Prots.arxiv.bbl}


\end{document}
